
% BEGIN preamble.tex
\documentclass[11pt,a4paper]{article}
\usepackage[T1]{fontenc}
\usepackage[utf8]{inputenc}
\usepackage{lmodern}
\usepackage{amsmath,amssymb,amsthm,mathtools}
\usepackage[margin=25mm,headheight=15pt]{geometry}
\usepackage{microtype,booktabs,longtable,array,enumitem}
\usepackage{xcolor,xurl}
\usepackage[colorlinks=true,linkcolor=blue!45!black,citecolor=blue!45!black,urlcolor=blue!45!black]{hyperref}
\usepackage{bookmark,fancyhdr}
\hypersetup{pdftitle={Quartic Digamma Bridges, Harmonic Tail Moments, and Complex Resolvent Powers},pdfauthor={Research continuation prepared with ChatGPT for Vladimir Reshetnikov and ProveIt},pdfsubject={Exact identities, complete endpoint normalizations, and reproducible verification}}
\numberwithin{equation}{section}
\newtheorem{theorem}{Theorem}[section]
\newtheorem{proposition}[theorem]{Proposition}
\newtheorem{lemma}[theorem]{Lemma}
\newtheorem{corollary}[theorem]{Corollary}
\theoremstyle{definition}
\newtheorem{definition}[theorem]{Definition}
\newtheorem{example}[theorem]{Example}
\theoremstyle{remark}
\newtheorem{remark}[theorem]{Remark}
\DeclareMathOperator{\Li}{Li}
\DeclareMathOperator{\Log}{Log}
\DeclareMathOperator{\FP}{FP}
\newcommand{\dd}{\mathop{}\!\mathrm{d}}
\newcommand{\snapshot}{\texttt{7bd45777a8a127e5dc69aff8887ca36a79a3059d}}
\newcommand{\shortsnapshot}{\texttt{7bd45777a8a1}}
\setlist{itemsep=4pt,topsep=5pt}
\setlength{\parskip}{3pt}
\setlength{\parindent}{15pt}
\setlength{\emergencystretch}{2em}
\allowdisplaybreaks[1]
\pagestyle{fancy}
\fancyhf{}
\fancyhead[L]{\small\scshape ProveIt research continuation}
\fancyhead[R]{\small 11 October 2026}
\fancyfoot[C]{\thepage}
\renewcommand{\headrulewidth}{0.3pt}
\setcounter{tocdepth}{2}


% END preamble.tex

\begin{document}
\begin{titlepage}
\thispagestyle{empty}
\vspace*{10mm}
\begin{center}
{\small\scshape Exact identities and spectral continuation\par}
\vspace{10mm}
{\LARGE\bfseries Quartic Digamma Bridges,\\[4pt]
Harmonic Tail Moments,\\[4pt]
and Complex Resolvent Powers\par}
\vspace{8mm}
{\large Complete proofs, endpoint corrections,\\
and normalized antiderivatives\par}
\vspace{10mm}
{Research continuation prepared with ChatGPT\\
for Vladimir Reshetnikov and the ProveIt project\par}
\vspace{3mm}
{11 October 2026 (UTC)\par}
\end{center}
\vspace{7mm}
\begin{abstract}
We resolve three concrete parts of the current ProveIt research agenda.
First, we evaluate the coordinate finite part of the fourth power of
the digamma function through the linear Laurent coefficient of the
strict depth-three harmonic zeta function, with every lower-depth and
endpoint correction explicit. A normally convergent global
partial-fraction identity supplies a normalized primitive and an
all-order polygamma hierarchy. Second, we evaluate every centered
nonpositive even moment of a product of two even harmonic tails by a
finite Bernoulli formula. For each pair of tail orders the entire
moment sequence belongs to one fixed span of ordinary zeta values and
products, and its convergent generating function reduces to finitely
many classical polylogarithms. The first mixed pair is generated
rationally by its first four moments. Third, we extend the cotangent
resolvent transform to arbitrary complex powers and complete its joint
Stieltjes germ at every integer power. The completion includes a finite
endpoint term that remains nonzero when the apparent pole is removable.
It gives all logarithmic powers, exact parameter derivatives, and
normalized primitives. An independent Gamma--Gauss generator evaluates
all ordinary polynomial moments. We separate these proofs from
computational diagnostics, record the source audit, and identify the
remaining arithmetic and higher-depth questions.
\end{abstract}
\vfill
\noindent\textbf{Fixed source baseline.} ProveIt revision
\shortsnapshot, including all eleven incoming research archives.
The precise source map and full commit identifier appear in
Section~\ref{intro:scope}. The Gaussian $S_6$ and revised $S_8$
conjectures remain open.
\end{titlepage}
{\footnotesize\setlength{\parskip}{0pt}\tableofcontents}
\clearpage

% BEGIN sections/introduction.tex
\section{Source baseline and mathematical contributions}
\label{intro:scope}

This article continues the polylogarithm manuscript and incoming
research reports in the ProveIt repository~\cite{ProveIt} at the fixed
revision
\begin{center}\snapshot.\end{center}
The source comparison uses the manuscript, the relevant archived
continuations, and all eleven incoming ZIP archives present at this
revision. The incoming directory is a research staging area: a claim
being present there is not evidence that it is proved or integrated
into a canonical manuscript chapter. We therefore trace each advance
to a named question and preserve the qualifications in its source.

The immediate starting points are the cubic harmonic
bridge~\cite{IncomingCubic}, the centered parity and rational-resolvent
theorems~\cite{IncomingParity}, and the higher-reflection
generators~\cite{IncomingReflection}. The strict ordered-Hurwitz
conventions come from the all-depth germ report~\cite{IncomingOrdered}.
The exact-relation obstruction and the incomplete larger relation
search~\cite{IncomingExact,IncomingIndependent} delimit what can
presently be claimed about the Gaussian polylogarithm conjectures.

\subsection{Three resolved research directions}

The main advances are the following. They are new relative to the
inspected source corpus; no exhaustive claim of priority across the
literature is intended.

\begin{enumerate}
\item \textbf{The explicitly requested quartic harmonic reduction.}
Question 6 of the cubic harmonic report asks for the finite reduction
at $k=4$, with the higher-depth data identified. We prove
Theorem~\ref{q4:main}, fixing the entire remainder in the quartic
Mittag--Leffler expansion and all integration constants. The answer
uses the linear Laurent coefficient $\theta$ of $Z_3(1+u,1,1)$,
the inherited coefficient $\eta$ of $Z_2(1+u,1)$, and ordinary
zeta and Stieltjes constants. Equation~\eqref{q4:single-coordinate}
eliminates $\eta$ by combining the fourth and third moments.
An absolutely convergent harmonic series, an elementary Mellin
kernel for $\theta$, a normalized primitive, and every argument
derivative accompany the evaluation.

\item \textbf{Every moment of every pair of even harmonic tails.}
Question R1 of the harmonic-parity report asks first for the mixed
pair of orders two and four, and then for a controlled algebra at
general orders. Theorems~\ref{tail:all-moments} and
\ref{tail:generator} solve the pure two-tail family for all even
$p,q\ge2$ and all moment orders. The finite formula retains no
unreduced double-zeta coordinate: all double zeta values have odd
weight and reduce by classical identities. The moment order does
not enlarge the specified arithmetic span. For $(p,q)=(2,4)$,
every moment is a rational combination of the first four. The
additional factors of $H_n$ and odd-order harmonic numbers mentioned
in R1 remain outside this theorem.

\item \textbf{Complex powers and full moving-endpoint completion.}
Question R4 of the same report asks for a branch, a joint
meromorphic generator, and the full resonant amplitudes for
$(\pi\cot\pi x-z)^{-\lambda}$. Theorems
\ref{power:eulerthm}, \ref{power:hurwitzthm}, and
\ref{power:completion} provide them. The result applies to every
nonreal $z$, every Stieltjes index and argument derivative, and
every integer power with arbitrary logarithmic powers. Away from
integer powers the uncompleted local formula applies directly.
The correction at negative integers is finite even when its
denominator is removable; omitting it changes the answer.
Section~\ref{moment:sec} adds a Gamma--Gauss generator for all
ordinary polynomial moments and a generalized-harmonic formula
at the imaginary base points.
\end{enumerate}

The status changes are deliberately narrower than the broadest
questions in the sources. The quartic calculation does not settle
every power $k\ge4$. The two-tail theorem does not treat four tails
or extra harmonic factors. A functional transform for arbitrary
complex powers need not reduce every one of its parameter values to
a finite list of familiar constants. Each theorem states exactly
what is finite, what is continued, and what arithmetic coordinate
remains.

\subsection{Representative identities}

Two compact consequences illustrate the range of the results.
With $x_n=n+1/2$, one obtains the ordinary absolutely convergent sum
\begin{equation}\label{intro:rational-example}
\begin{aligned}
\sum_{n=0}^{\infty}\bigg\{&
\frac{1344x_n^6+2000x_n^4+412x_n^2+7}{6}
                  \psi'(n+1)\psi'''(n+1)\\
&\hspace{17mm}-448x_n^2-\frac{1216}{3}\bigg\}=-88.
\end{aligned}
\end{equation}
The polynomial subtraction is part of the summand; it is not a
separate divergent series. The proof is exact elimination among
four independently specified centered moments, as detailed in
\eqref{tail:24-eliminations}--\eqref{tail:rational-polygamma}.

At a negative half power of the resolvent, an ordinary real integral
has the elementary zeta/logarithm evaluation
\begin{equation}\label{intro:half-example}
\int_0^1\left(x^2-x+\frac16\right)
       \frac{\cos(\pi x/2)}{\sqrt{\sin(\pi x)}}\,\dd x
       =\frac1{12}-\frac{\log^22}{\pi^2}.
\end{equation}
Its entire sequence of Bernoulli analogues follows from one Gamma
quotient, rather than separate evaluations of successive moments.

These ordinary identities complement the finite-part identity for
$\psi^4$. A finite part is appropriate in that case because the
unsubtracted integral diverges at zero; it is not required for
\eqref{intro:rational-example} or \eqref{intro:half-example}.

\subsection{Classical inputs and proof architecture}

The classical ingredients are digamma recurrence and reflection,
Hurwitz Fourier expansions, Mellin transforms, Bernoulli
polynomials, Euler sums, and beta/hypergeometric integral formulas
\cite{DLMF,BBG,EspinosaMoll}. Harmonic-zeta and Stieltjes
coordinates have their own established literature; Coppo's
harmonic-zeta identities provide a recent point of comparison
\cite{Coppo}. We use these sources for specified formulas, not to
infer any unproved constant reduction.

Three proof mechanisms perform different jobs. A two-point contour
fixes the entire function that remains after matching the digamma
principal parts. Exact finite summation by parts identifies the
constant at infinity in a harmonic-tail moment without rearranging
a divergent double sum. Finally, local monomial subtraction in two
complex parameters distinguishes a meromorphic spectral value
from the coordinate finite part of its specialized integrand.
The last distinction is essential for differentiation and
antidifferentiation at resonances.

The proofs do not depend on a numerical relation search. The
accompanying scripts check finite algebra independently and compare
analytically different numerical representations. Explicit
truncation estimates support these comparisons, but the recorded
floating-point runs do not constitute interval certificates or
formal proof-assistant proofs.


% END sections/introduction.tex


% BEGIN sections/conventions.tex
\section{Conventions and finite-part normalization}
\label{conv:section}

We use $\psi=\Gamma'/\Gamma$, $\psi^{(0)}=\psi$, and
$H_n^{(r)}=\sum_{j=1}^n j^{-r}$, with $H_n=H_n^{(1)}$ and
$H_0^{(r)}=0$. A prime on $\zeta(s)$ denotes differentiation
in its order $s$; a superscript $[k]$ on $\Li_s$ also denotes
order differentiation. Derivatives of $\gamma_m(x)$ are with
respect to its argument $x$.

The generalized Stieltjes convention is
\begin{equation}\label{conv:stieltjes}
\zeta(s,x)=\frac1{s-1}+
\sum_{m\ge0}\frac{(-1)^m\gamma_m(x)}{m!}(s-1)^m.
\end{equation}
Thus $\gamma_0(x)=-\psi(x)$ and
$\gamma_m=\gamma_m(1)$. After setting $s=1-u$, the pole-removed
germ is
\[
\zeta(1-u,x)+\frac1u
=\sum_{m\ge0}\gamma_m(x)\frac{u^m}{m!}.
\]
This positive-sign $u$ expansion must not be confused with the
alternating-sign expansion of $\zeta(1+u)$.

Strict multiple zeta functions have decreasing indices:
\[
Z_r(s_1,\ldots,s_r)=
\sum_{n_1>\cdots>n_r\ge1}\prod_{j=1}^r n_j^{-s_j}.
\]
We also write $\zeta(a,b)=Z_2(a,b)$ when the sum converges.
The coefficients $\eta$ and $\theta$ in Section~\ref{q4:section}
are coefficients of $u^1$, with no hidden factorial or sign.
Bernoulli numbers have $B_1=-1/2$, and Bernoulli polynomials use
\[
\frac{te^{xt}}{e^t-1}=\sum_{n\ge0}B_n(x)\frac{t^n}{n!}.
\]

\begin{definition}[Coordinate finite part]
Suppose an integrand on $(0,1)$ has, to every required finite order,
endpoint expansions in complex powers and powers of logarithms,
with an integrable remainder after finitely many subtractions.
Its coordinate finite part is the constant coefficient of the
cutoff integral in the endpoint coordinates $x$ and $1-x$.
One may introduce independent cutoffs at the two ends; the
constant is the coefficient of both cutoff powers zero and both
logarithm powers zero. A common cutoff gives the same constant
for the separate endpoint expansions used here.
\end{definition}

For example,
\[
\FP\int_0^1x^{-j}\,\dd x=-\frac1{j-1}\quad(j\ge2),
\qquad
\FP\int_0^1\frac{\log^k x}{x}\,\dd x=0\quad(k\ge0).
\]
At a general nonzero exponent $\alpha$ the regularized monomial
integral is $1/\alpha$; its parameter derivatives are the
corresponding logarithmic integrals. At $\alpha=0$ the coordinate
constant is zero, whereas the meromorphic function has a pole.
This elementary difference becomes a finite contact term when a
coefficient also vanishes at the same parameter.

A change of endpoint coordinate can change a finite part. Every
identity in this article fixes the coordinate before taking the
constant, differentiating parameters, or forming a primitive.
In particular, the spectral values $D_{p,q}(-2r)$ are initially
defined by meromorphic continuation, and their equality to
ordinary subtracted sums is proved in
Proposition~\ref{tail:ordinary}; it is not an implicit choice
of summation method.

All primitives are normalized. The quartic primitive has zero
constant term at its singular origin. Primitives in the nonreal
resolvent parameter are differences at a specified base point
within one half-plane. Complex powers of the resolvent use the
continuous logarithm specified in
Section~\ref{power:sec}; no branch is inferred from an
unqualified power of a complex number.


% END sections/conventions.tex


% BEGIN sections/quartic_bridge.tex
% Research section prepared against ProveIt 7bd45777a8a127e5dc69aff8887ca36a79a3059d.
% This fragment uses only standard amsmath/amsthm commands.  Labels begin q4:.
% Suggested source citations are recorded in SOURCE_NOTES.md beside this file.

\section{The quartic digamma finite part and a depth-three harmonic coefficient}
\label{q4:section}

The cubic harmonic report~\cite{IncomingCubic} leaves a specific
next question: carry out the
power-to-harmonic reduction for the fourth power of the digamma function,
with the entire remainder and every endpoint constant fixed.  The theorem
below completes that reduction.  The retained depth-three coefficient is
defined explicitly; a finite evaluation of it in ordinary zeta and
Stieltjes constants is a separate arithmetic question.

Throughout this section, write $\gamma=\gamma_0$, and use
\[
 \zeta(1+u)=\frac1u+\sum_{j\geq0}\frac{(-1)^j\gamma_j}{j!}u^j.
\]
For strict multiple zeta functions the larger index is first:
\[
 Z_r(s_1,\ldots,s_r)
 =\sum_{n_1>\cdots>n_r\geq1}n_1^{-s_1}\cdots n_r^{-s_r}.
\]
Only a single spectral order varies in the following definitions:
\begin{align}
 Z_2(1+u,1)
 &=u^{-2}+\gamma u^{-1}+\frac{\gamma^2-\zeta(2)}2
       +\eta u+O(u^2),\label{q4:eta}\\
 Z_3(1+u,1,1)
 &=u^{-3}+\gamma u^{-2}
       +\frac{\gamma^2-\zeta(2)}{2u}
       +\frac{\gamma^3-3\gamma\zeta(2)+2\zeta(3)}6
       +\theta u+O(u^2).\label{q4:theta}
\end{align}
Thus $\theta$ is the linear Laurent coefficient, with no factorial or sign
included.  These conventions agree with the incoming ordered-Hurwitz and
cubic harmonic reports~\cite{IncomingOrdered,IncomingCubic}.
The expansions can also be established directly
from the finite elementary symmetric sums used below.

Define $\operatorname{FP}\int_0^1$ as the constant coefficient of
$\int_\varepsilon^1$ in the coordinate $\varepsilon$ itself.  In particular,
$\operatorname{FP}\int_0^1x^{-j}\,dx=-1/(j-1)$ for $j\geq2$, whereas
$\operatorname{FP}\int_0^1x^{-1}\,dx=0$.

\begin{theorem}[Quartic harmonic bridge]\label{q4:main}
Let
\[
 A_4=12\zeta(3)-12\gamma\zeta(2)+18\zeta(4)+4\zeta'(3).
\]
Then
\begin{equation}\label{q4:bridge}
 \boxed{\displaystyle
 Q_4:=\operatorname{FP}\int_0^1\psi(x)^4\,dx
 =A_4-24\theta+24\gamma\eta
        +12\bigl(\gamma^2-\zeta(2)\bigr)\gamma_1.}
\end{equation}
Equivalently, an absolutely convergent arithmetic representation is
\begin{equation}\label{q4:arithmetic}
 \boxed{\begin{aligned}
 Q_4={}&A_4+12\gamma_3-24\zeta(2)\gamma_1\\
 &+12\sum_{n=1}^{\infty}
 \frac{\bigl(\psi(n)^2+\psi'(n)-\log^2n\bigr)\log n}{n}.
 \end{aligned}}
\end{equation}
In generalized harmonic notation the numerator inside the parentheses is
\[
 (H_{n-1}-\gamma)^2+\zeta(2)-H_{n-1}^{(2)}-\log^2n.
\]
Using the already proved cubic bridge, one may isolate the single
depth-three coefficient in
\begin{align}\label{q4:single-coordinate}
 \operatorname{FP}\int_0^1\bigl\{\psi(x)^4+4\gamma\psi(x)^3\bigr\}\,dx
 ={}&12\zeta(3)+18\zeta(4)+4\zeta'(3)\notag\\
 &-24\theta-12\bigl(\gamma^2+\zeta(2)\bigr)\gamma_1.
\end{align}
No arithmetic independence or ordinary-constant evaluation of $\theta$ is
asserted by these formulas.
\end{theorem}

\subsection{A global quartic partial-fraction identity}

For $n\geq0$ put
\begin{equation}\label{q4:coefficients}
 \begin{split}
 b_n&=H_n-\gamma,\qquad
 c_n=\zeta(2)+H_n^{(2)},\qquad
 d_n=H_n^{(3)}-\zeta(3),\\
 q_n&=6b_n^2-4c_n,\qquad
 a_n=-b_n^3+3b_nc_n-d_n,
 \end{split}
\end{equation}
with $H_0^{(r)}=0$.  The quartic principal part at $-n$ is
\begin{equation}\label{q4:Pn}
 P_n(z)=\frac1{(z+n)^4}-\frac{4b_n}{(z+n)^3}
             +\frac{q_n}{(z+n)^2}+\frac{4a_n}{z+n}.
\end{equation}
The constant coefficient of $\psi(z)^4$ at $z=0$ is
\begin{equation}\label{q4:kappa}
 \kappa_4=\gamma^4-12\gamma^2\zeta(2)+6\zeta(2)^2
                        +12\gamma\zeta(3)-4\zeta(4).
\end{equation}

\begin{lemma}[The entire remainder is fixed]\label{q4:ML}
For $z\notin\{0,-1,-2,\ldots\}$,
\begin{equation}\label{q4:MLformula}
 \psi(z)^4=P_0(z)+\kappa_4+
                   \sum_{n=1}^{\infty}\bigl(P_n(z)-P_n(0)\bigr).
\end{equation}
The series converges normally on every compact set avoiding the poles.
\end{lemma}

\begin{proof}
The digamma recurrence and the Taylor expansion at zero
\cite[Sections~5.5, 5.7]{DLMF} give
\[
 \psi(-n+\varepsilon)=-\varepsilon^{-1}+b_n+c_n\varepsilon
           +d_n\varepsilon^2
           +\bigl(\zeta(4)+H_n^{(4)}\bigr)\varepsilon^3
           +O(\varepsilon^4).
\]
Taking its fourth power proves \eqref{q4:Pn}; its constant term at zero
is \eqref{q4:kappa}.  Matching principal parts by itself would leave an
unknown entire function, so one additional argument is necessary.

Let $f(z)=\psi(z)^4-P_0(z)$, which is regular at zero.  For fixed $z$,
integrate
\[
 \frac{z f(w)}{w(w-z)}
\]
over $|w|=N+1/2$, with $N$ sufficiently large.  On these circles,
$f(w)=O(\log^4N)$ uniformly.  In the right half-plane this follows from
the usual digamma asymptotic~\cite[Section~5.11]{DLMF}.
In the left half-plane use
$\psi(w)=\psi(1-w)-\pi\cot(\pi w)$.  The cotangent is bounded uniformly
on the chosen circles: for $|\operatorname{Im}w|\geq1/4$ its exponential
formula suffices, and for the remaining arcs the real part remains a
fixed distance from every integer.  Hence the contour integral is
$O(\log^4N/N)$ and tends to zero.

Its residues at $z$ and $0$ are $f(z)$ and $-f(0)$, and the residue at
$-n$ is $-(P_n(z)-P_n(0))$.  The residue theorem proves
\eqref{q4:MLformula}.  Finally,
$b_n=O(\log n)$, $q_n=O(\log^2n)$ and $a_n=O(\log^3n)$, while
$(n+z)^{-1}-n^{-1}=O(n^{-2})$ uniformly on a fixed compact set.
This proves normal convergence as well as the identity.
\end{proof}

\subsection{Integration and finite summation by parts}

Integrating the normally convergent remainder in \eqref{q4:MLformula}
term by term gives
\begin{align}\label{q4:integrated}
 Q_4={}&-\frac13-2\gamma-q_0+\kappa_4\notag\\
 &+\sum_{n\geq1}\left\{
   \frac{n^{-3}-(n+1)^{-3}}3-n^{-4}
   -2b_n\bigl(n^{-2}-(n+1)^{-2}\bigr)+4b_nn^{-3}\right.\notag\\
 &\hspace{34mm}\left.
   +q_n\bigl(n^{-1}-(n+1)^{-1}-n^{-2}\bigr)
   +4a_n\left(\log\left(1+\frac1n\right)-\frac1n\right)
  \right\}.
\end{align}
The finite parts $-1/3$, $-2\gamma$ and $-q_0$ are retained explicitly.
The simple-pole finite part at the origin is zero.

We use the following classical convergent Euler sums
\cite[Equations~25.16.8 and 25.16.13]{DLMF}:
\begin{align*}
 \sum_{n\geq1}\frac{H_n}{n^2}&=2\zeta(3),&
 \sum_{n\geq1}\frac{H_n}{n^3}&=\frac54\zeta(4),\\
 \sum_{n\geq1}\frac{H_n^2}{n^2}&=\frac{17}4\zeta(4),&
 \sum_{n\geq1}\frac{H_n^{(2)}}{n^2}
     &=\frac{\zeta(2)^2+\zeta(4)}2.
\end{align*}
The first three can be taken from the standard Euler-sum formulas; the
fourth follows by splitting the product $\zeta(2)^2$ into its two strict
triangles and its diagonal.  Also $\zeta(2)^2=5\zeta(4)/2$.
The exact differences
\begin{equation}\label{q4:differences}
 b_n-b_{n-1}=\frac1n,\qquad
 q_n-q_{n-1}=\frac{12b_{n-1}}n+\frac2{n^2},\qquad
 a_n-a_{n-1}=\frac{3(c_{n-1}-b_{n-1}^2)}n+\frac1{n^3}
\end{equation}
then reduce all terms in \eqref{q4:integrated} except the logarithmic
ones to
\begin{equation}\label{q4:base}
 Q_4=B_4+4\sum_{n\geq1}a_n
          \left(\log\left(1+\frac1n\right)-\frac1n\right),
\end{equation}
where
\[
 B_4=\gamma^4-18\gamma^2\zeta(2)-12\gamma\zeta(2)
             +32\gamma\zeta(3)+12\zeta(3)+\frac{13}2\zeta(4).
\]
For a directly checkable bookkeeping formula, the lower-pole terms before
substitution of the Euler sums are
\[
 \kappa_4-\zeta(4)
 +4\sum_{n\geq1}\frac{b_n}{n^3}
 +12\sum_{n\geq1}\frac{b_{n-1}}{n^2}
 -6\sum_{n\geq1}\frac{b_n^2}{n^2}
 +4\sum_{n\geq1}\frac{c_n}{n^2}.
\]
This equals $B_4$ exactly.

The residue difference in \eqref{q4:differences} is the decisive step.
For a finite integer $N\geq1$, define
\[
 M_N=\sum_{n=1}^N
             \frac{(b_{n-1}^2-c_{n-1})\log n}{n}.
\]
Finite summation by parts gives the exact equality
\begin{equation}\label{q4:logsum}
 \sum_{n=1}^N a_n\log\left(1+\frac1n\right)
       =a_N\log(N+1)+3M_N-\sum_{n=1}^N\frac{\log n}{n^3}.
\end{equation}
To evaluate the other finite sum, let $e_j(N)$ be the elementary symmetric
sum of degree $j$ in $1,1/2,\ldots,1/N$, and write
\[
 H_N(r_1,\ldots,r_d)
  =\sum_{N\geq n_1>\cdots>n_d\geq1}
         n_1^{-r_1}\cdots n_d^{-r_d}.
\]
The polynomial expression for the residue is
\begin{equation}\label{q4:elementary}
 a_n=-6e_3(n)+6\gamma e_2(n)+3(\zeta(2)-\gamma^2)H_n
           +\gamma^3-3\gamma\zeta(2)+\zeta(3)+H_n^{(3)}.
\end{equation}
Use $e_j(n)=e_j(n-1)+e_{j-1}(n-1)/n$ in each sum, and split the
product $H_NH_N^{(3)}$ into strict triangles.  This proves, without
any infinite rearrangement,
\begin{align}\label{q4:finitesum}
 \sum_{n=1}^N\frac{a_n}{n}
 ={}&-6e_4(N)+6\gamma e_3(N)+3(\zeta(2)-\gamma^2)e_2(N)\notag\\
 &+(\gamma^3-3\gamma\zeta(2)+\zeta(3))H_N+H_NH_N^{(3)}\notag\\
 &-6H_N(2,1,1)+6\gamma H_N(2,1)
       +3(\zeta(2)-\gamma^2)H_N^{(2)}-H_N(3,1).
\end{align}
The only multiple-zeta limits needed here are
\[
 Z_3(2,1,1)=\zeta(4),\qquad
 Z_2(2,1)=\zeta(3),\qquad
 Z_2(3,1)=\frac14\zeta(4).
\]
For example,
\[
 Z_3(2,1,1)=\frac16\int_0^1\frac{[-\log(1-t)]^3}{t}\,dt
           =\frac16\int_0^1\frac{(-\log t)^3}{1-t}\,dt
           =\zeta(4),
\]
and the double-zeta values follow from the displayed Euler sums.

The elementary symmetric formulas are
\[
 \begin{split}
 e_2(N)&=\frac{H_N^2-H_N^{(2)}}2,\\
 e_3(N)&=\frac{H_N^3-3H_NH_N^{(2)}+2H_N^{(3)}}6,\\
 e_4(N)&=\frac{H_N^4-6H_N^2H_N^{(2)}+3(H_N^{(2)})^2
                +8H_NH_N^{(3)}-6H_N^{(4)}}{24}.
 \end{split}
\]
Substitute these into \eqref{q4:finitesum}, using
$H_N=\log N+\gamma+O(N^{-1})$ and
$H_N^{(r)}=\zeta(r)+O(N^{1-r})$ for $r\geq2$.
All discarded terms are $o(1)$, including their polynomial logarithmic
factors.  The result, including its full constant term, is
\begin{equation}\label{q4:asymptoticsum}
 \sum_{n=1}^N\frac{a_n}{n}
 =-\frac{\log^4N}{4}+3\zeta(2)\log^2N+C_4+o(1),
\end{equation}
where
\begin{equation}\label{q4:C4}
 C_4=\frac{\gamma^4}{4}-\frac92\gamma^2\zeta(2)
                    +8\gamma\zeta(3)-\frac{23}8\zeta(4).
\end{equation}
Also
\[
 a_N\log(N+1)=-\log^4N+6\zeta(2)\log^2N+o(1).
\]
Equations \eqref{q4:base}--\eqref{q4:asymptoticsum} now yield
\begin{equation}\label{q4:limit}
 Q_4=A_4+12\lim_{N\to\infty}
       \left\{M_N-\frac{\log^4N}{4}+\zeta(2)\log^2N\right\}.
\end{equation}
Indeed $B_4-4C_4=12\zeta(3)-12\gamma\zeta(2)+18\zeta(4)$,
and $-\sum\log n/n^3=\zeta'(3)$.  This is where an omitted endpoint
constant or a sign error in the logarithmic sum would alter the theorem.

\subsection{Identification of the Laurent coefficient}

Since
\[
 b_{n-1}^2-c_{n-1}
   =\psi(n)^2+\psi'(n)-2\zeta(2),
\]
and $\psi(n)=\log n+O(n^{-1})$, $\psi'(n)=O(n^{-1})$, the series
in \eqref{q4:arithmetic} is absolutely convergent, with summands
$O(\log^2n/n^2)$.  The standard Stieltjes limits give
\begin{align}\label{q4:Mconstant}
 \lim_{N\to\infty}
       \left\{M_N-\frac{\log^4N}{4}+\zeta(2)\log^2N\right\}
  ={}&\gamma_3-2\zeta(2)\gamma_1\notag\\
 &+\sum_{n\geq1}
       \frac{(\psi(n)^2+\psi'(n)-\log^2n)\log n}{n}.
\end{align}
Together with \eqref{q4:limit}, this proves \eqref{q4:arithmetic}.

For the precise spectral interpretation, define the holomorphic function
\[
 D_4(s)=\sum_{n\geq1}
       \frac{\psi(n)^2+\psi'(n)-\log^2n}{n^s},
       \qquad\operatorname{Re}s>0.
\]
The series and every fixed derivative converge normally on compact
subsets of this half-plane.  In $\operatorname{Re}s>1$ the finite identity
$H_{n-1}^2-H_{n-1}^{(2)}=2e_2(n-1)$ gives
\begin{equation}\label{q4:Didentity}
 D_4(s)=2Z_3(s,1,1)-2\gamma Z_2(s,1)
                  +(\gamma^2+\zeta(2))\zeta(s)-\zeta''(s).
\end{equation}
Analytic continuation carries the identity to the germ at $s=1$.
Comparing its linear coefficients, with \eqref{q4:eta} and
\eqref{q4:theta}, yields
\[
 D_4'(1)=2\theta-2\gamma\eta
                    -(\gamma^2+\zeta(2))\gamma_1+\gamma_3.
\]
The convergent sum in \eqref{q4:arithmetic} is $-D_4'(1)$.
Substitution proves \eqref{q4:bridge}.  Finally, the incoming cubic
report~\cite{IncomingCubic}
proves
\[
 Q_3=\operatorname{FP}\int_0^1\psi(x)^3\,dx
    =3\zeta(2)-6\eta-6\gamma\gamma_1.
\]
Adding $4\gamma Q_3$ gives \eqref{q4:single-coordinate}.

\paragraph{A direct existence check for the depth-three Laurent germ.}
There is also a useful elementary Mellin representation:
\begin{equation}\label{q4:heightone-mellin}
 Z_3(s,1,1)=\frac1{2\Gamma(s)}\int_0^\infty
     \frac{t^{s-1}\log^2(1-e^{-t})}{e^t-1}\,dt,
     \qquad\operatorname{Re}s>1.
\end{equation}
Indeed the generating series of $e_2(n-1)$ is
$\log^2(1-z)/(2(1/z-1))$.  Mellin integration is absolutely convergent
in the stated domain.  Near zero the kernel is
$\log^2t/(2t)+O(\log^2t)$, so subtracting that one local term gives
a holomorphic remainder $R_3(u)$ in
\[
 \Gamma(1+u)Z_3(1+u,1,1)=u^{-3}+R_3(u).
\]
The identity
$\frac{d}{dt}\log^3(1-e^{-t})=3\log^2(1-e^{-t})/(e^t-1)$
shows $R_3(0)=0$.  The inverse-Gamma expansion therefore gives exactly
the principal and constant coefficients in \eqref{q4:theta}.  In
particular,
\begin{align}\label{q4:theta-kernel}
 \theta={}&\frac{\gamma^4-6\gamma^2\zeta(2)+3\zeta(2)^2
                       +8\gamma\zeta(3)-6\zeta(4)}{24}\notag\\
 &+\frac12\int_0^1\log t
       \left\{\frac{\log^2(1-e^{-t})}{e^t-1}
                        -\frac{\log^2t}{t}\right\}\,dt\notag\\
 &+\frac12\int_1^\infty
       \frac{\log t\,\log^2(1-e^{-t})}{e^t-1}\,dt.
\end{align}
Both integrals in this display converge ordinarily.  This fixes the
retained coordinate by an elementary, one-dimensional kernel as well as
by the strict Laurent definition; it does not remove its arithmetic
content by renaming it.

\subsection{A normalized quartic antiderivative and derivative identities}

For $j\geq2$, define
\[
 F_j(n,z)=\frac{(n+z)^{1-j}-n^{1-j}}{1-j}-\frac{z}{n^j},
 \qquad
 F_1(n,z)=\log\left(1+\frac zn\right)-\frac zn.
\]
Use the logarithm real on the positive real axis, continued to
$\operatorname{Re}z>0$.  Put
$p_{4,n}=1$, $p_{3,n}=-4b_n$, $p_{2,n}=q_n$, $p_{1,n}=4a_n$.

\begin{corollary}[Normalized primitive]\label{q4:primitive}
On $\operatorname{Re}z>0$, the function
\begin{align}\label{q4:primitiveformula}
 \mathcal P_4(z)={}&-\frac1{3z^3}-\frac{2\gamma}{z^2}
       -\frac{6\gamma^2-4\zeta(2)}z
       +4(\gamma^3-3\gamma\zeta(2)+\zeta(3))\log z
       +\kappa_4z\notag\\
 &+\sum_{n=1}^\infty\sum_{j=1}^4p_{j,n}F_j(n,z)
\end{align}
satisfies $\mathcal P_4'(z)=\psi(z)^4$.  Its expansion at the singular
origin has constant term zero.  Consequently, for every real $b>0$,
\[
 \operatorname{FP}\int_0^b\psi(x)^4\,dx=\mathcal P_4(b),
 \qquad \mathcal P_4(1)=Q_4,
\]
and, for $0<a<b$, the ordinary integral is
$\int_a^b\psi(x)^4\,dx=\mathcal P_4(b)-\mathcal P_4(a)$.
\end{corollary}

\begin{proof}
Each $F_j$ is the primitive vanishing at zero of
$(n+z)^{-j}-n^{-j}$.  Uniformly on a fixed compact set,
$F_j(n,z)=O(n^{-j-1})$.  The coefficient growth established in
Lemma~\ref{q4:ML} therefore proves normal convergence of the series and
its differentiated series.  Differentiation gives
\eqref{q4:MLformula}.  The same estimates give uniform convergence on
each disk $|z|\leq r<1$, including the origin, where every summand
vanishes.  Thus the displayed singular terms fix the constant
coefficient to be zero.
\end{proof}

\begin{corollary}[Quartic polygamma hierarchy]\label{q4:derivatives}
For every integer $r\geq1$, away from the poles of $\psi$,
\begin{align}\label{q4:allpolygamma}
 &\sum_{r_1+r_2+r_3+r_4=r}
       \frac{r!}{r_1!r_2!r_3!r_4!}
       \prod_{j=1}^4\psi^{(r_j)}(z)\notag\\
 &\quad=(-1)^rr!\sum_{n=0}^{\infty}
       \left\{\frac{\binom{r+3}{3}}{(n+z)^{r+4}}
       -\frac{4b_n\binom{r+2}{2}}{(n+z)^{r+3}}
       +\frac{(r+1)q_n}{(n+z)^{r+2}}
       +\frac{4a_n}{(n+z)^{r+1}}\right\}.
\end{align}
The series converges normally.  At $z=1$ its value is a finite polynomial
in $\gamma,\zeta(2),\ldots,\zeta(r+1)$.
\end{corollary}
\begin{proof}
Differentiate \eqref{q4:MLformula} $r$ times.  Its constant subtractions
disappear, and the resulting terms are bounded by a constant times
$n^{-r-1}(1+\log^3n)$ on each compact set avoiding the poles.
The product rule gives the left side.  The last assertion follows from
$\psi(1)=-\gamma$ and
$\psi^{(k)}(1)=(-1)^{k+1}k!\zeta(k+1)$ for $k\geq1$.
\end{proof}

\subsection{Independent verification and a reproducible tail formula}

The supplied program checks 44 exact algebraic identities, including
all principal coefficients, the residue differences, the finite
summation identity for $1\leq N\leq25$, the endpoint-constant
cancellation, Laurent coefficient extraction, and the primitive
derivatives.  Its separate numerical calculations use 100 decimal
digits and give
\[
 Q_4=13.510120376015827411606240712830448441932798533301918600986096436\ldots.
\]
One evaluator uses the arithmetic series with a Binet tail.  The other
integrates the convergent Taylor series of $(x\psi(x))^4$ near zero and
uses ordinary quadrature away from zero; it does not use the harmonic
bridge or the Mittag--Leffler expansion.  Two choices of the matching
point, $1/8$ and $1/16$, agree to the displayed precision.  The finest
arithmetic result differs from these evaluations by less than
$1.40\cdot10^{-83}$.

Here is the explicit analytic tail control used in the arithmetic
calculation.  For $K\geq0$ set
\[
 a_K(x)=-\frac1{2x}-\sum_{k=1}^K\frac{B_{2k}}{2k\,x^{2k}},
 \qquad
 p_K(x)=\frac1x+\frac1{2x^2}
                      +\sum_{k=1}^K\frac{B_{2k}}{x^{2k+1}}.
\]
Binet's integral~\cite[Equation~5.9.15]{DLMF} and the finite
geometric expansion give, for $x>0$,
\[
 |\psi(x)-\log x-a_K(x)|\leq E_Kx^{-2K-2},\qquad
 |\psi'(x)-p_K(x)|\leq G_Kx^{-2K-3},
\]
where $G_K=|B_{2K+2}|$ and $E_K=G_K/(2K+2)$.
For the differentiated bound, expand $(1+y)^{-2}$ under the differentiated
Binet integral; its Taylor remainder after $K$ terms has magnitude at
most $(K+1)y^K$ for $y\geq0$.

For $N\geq2$ put
$U_{N,K}=1/(2N)+\sum_{k=1}^K|B_{2k}|/(2kN^{2k})$.
Replace $\psi(n)^2+\psi'(n)-\log^2n$ in the tail $n\geq N$ by
\[
 2\log n\,a_K(n)+a_K(n)^2+p_K(n).
\]
This finite polynomial in $n^{-1}$ and $\log n$ sums exactly by the first
two spectral derivatives of $\zeta(s,N)$.  The error in the series
itself is bounded by
\begin{align}\label{q4:tailbound}
 &2E_K\zeta''(2K+3,N)
    -2U_{N,K}E_K\zeta'(2K+3,N)\notag\\
 &\qquad{}-E_K^2\zeta'(4K+5,N)
    -G_K\zeta'(2K+4,N).
\end{align}
To obtain the error in $Q_4$, multiply this bound by $12$.
Indeed, if $R=\psi-\log-a_K$, then
\[
 |\psi^2+\psi'-(\log+a_K)^2-p_K|
 \leq2(\log n+U_{N,K})E_Kn^{-2K-2}
          +E_K^2n^{-4K-4}+G_Kn^{-2K-3},
\]
and summing after multiplication by $\log n/n$ gives
\eqref{q4:tailbound}.  With $N=64$, $K=34$, the bound for $Q_4$ is
below $1.88\cdot10^{-83}$.  The implementation uses ordinary mpmath
arithmetic, so this analytic truncation bound does not include an
interval enclosure of rounding error; the numerical evidence is
reported with that qualification.

\paragraph{Scope of the advance.}
The Mittag--Leffler method, harmonic Laurent coefficients, the Mellin
representation, and the Euler sums are classical.  The incoming cubic
report supplies the immediate proof strategy and explicitly asks for
the quartic reduction.  The new result relative to the inspected corpus
is \eqref{q4:bridge}, including its exact finite correction, convergent
arithmetic form and normalized primitive.  The claim is a solution of
that stated quartic reduction problem, not a proof that $\theta$ is
irreducible or a claim of global priority for every equivalent formula.

\paragraph{Further questions.}
The contour argument extends to each fixed power $\psi^k$, but the finite
arithmetic reduction has only been carried out here through $k=4$.
The next case asks which strict harmonic Laurent coefficients occur in
$\operatorname{FP}\int_0^1\psi(x)^5\,dx$ and whether some canonical
polynomial in $\psi$ cancels all lower-depth coefficients.  A second
question is to combine the new quartic bridge with the cubic logarithmic
sine moment supplied by the existing beta generator and determine the
smallest remaining spectral coordinate space.  A third is to obtain a
common-Hurwitz-shift version with all boundary terms fixed; simply
replacing $\gamma_j$ by $\gamma_j(a)$ in \eqref{q4:bridge} is not
justified by this proof.

% END sections/quartic_bridge.tex


% BEGIN sections/mixed_even_tails.tex
\section{Every centered moment of two even harmonic tails}
\label{tail:section}

The latest harmonic-parity report asks for exact centered moments of
\((\zeta(2i)-H_n^{(2i)})(\zeta(2j)-H_n^{(2j)})\), beginning with
\((i,j)=(1,2)\). Its all-degree evaluation covers only the trigamma
square. The following results answer the question for every pair of
even orders and every nonpositive even spectral argument. They also
give a convergent generating function in classical polylogarithms.
Products of four or more tails, and products with additional powers
of \(H_n\), are outside the theorem.

\subsection{Definitions, continuation, and an ordinary convergent sum}

Throughout this section, \(p,q\ge2\) are even integers, \(W=p+q\), and
\[
 A_p(n)=\zeta(p)-H_n^{(p)}=\zeta(p,n+1)
       =\frac{\psi^{(p-1)}(n+1)}{(p-1)!},
 \qquad x_n=n+\tfrac12.
\]
The last equality uses that \(p\) is even. Define
\begin{equation}\label{tail:D}
 D_{p,q}(s)=\sum_{n=0}^{\infty}\frac{A_p(n)A_q(n)}{x_n^s},
 \qquad \Re s>3-p-q.
\end{equation}
The defining series converges normally in this half-plane. A value
outside it always refers to meromorphic continuation unless it is
written as an explicitly subtracted convergent sum.

Set
\begin{equation}\label{tail:asymptotic-coefficients}
 c_{p,j}=\frac1{p-1}\binom{p+2j-2}{p-2}B_{2j}(1/2),
 \qquad
 k_{p,q,\ell}=\sum_{j=0}^{\ell}c_{p,j}c_{q,\ell-j}.
\end{equation}
Thus \(c_{p,0}=1/(p-1)\). The standard Hurwitz-zeta expansion gives
\begin{equation}\label{tail:centered-expansion}
 A_p(n)\sim\sum_{j\ge0}c_{p,j}x_n^{1-p-2j},
 \qquad
 A_p(n)A_q(n)\sim\sum_{\ell\ge0}
                  k_{p,q,\ell}x_n^{2-W-2\ell}.
\end{equation}
Every exponent in the product is even. The expansions and their
fixed-order remainder estimates follow either from the classical
Hurwitz formulas or directly from the Laplace integral used below
\cite{DLMF}.

\begin{proposition}[Continuation and ordinary summation]
\label{tail:ordinary}
The function \(D_{p,q}\) continues meromorphically to the plane. Its
only possible poles are simple and lie at
\[
 3-p-q,\ 1-p-q,\ -1-p-q,\ldots.
\]
In particular it is regular at every \(s=-2r\), \(r\ge0\).
Let
\[
 K_r=\max\{0,r-W/2+2\}.
\]
Then
\begin{equation}\label{tail:ordinary-formula}
 D_{p,q}(-2r)=\sum_{n=0}^{\infty}
 \left\{x_n^{2r}A_p(n)A_q(n)
       -\sum_{\ell=0}^{K_r-1}
         k_{p,q,\ell}x_n^{2r+2-W-2\ell}\right\}.
\end{equation}
An empty inner sum is zero. Every bracket is \(O(n^{-2})\) or
smaller, so the right side is an ordinary absolutely convergent sum.
\end{proposition}

\begin{proof}
Subtract any fixed number of terms from the product expansion in
\eqref{tail:centered-expansion}. The remainder series extends to a
larger left half-plane; the removed terms sum to
\(k_{p,q,\ell}\zeta(s+W-2+2\ell,1/2)\). Increasing the number of
subtractions proves the continuation and pole list. At \(s=-2r\),
the nondecaying terms are exactly those with \(0\le\ell<K_r\).
Their continued sums vanish, because
\(\zeta(-2m,1/2)=0\) for every integer \(m\ge0\). The remaining
sum converges absolutely and gives \eqref{tail:ordinary-formula}.
\end{proof}

The half-shift is essential to this argument: it makes both the
asymptotic powers and the polynomial subtraction compatible with the
trivial zeros of the same Hurwitz zeta function.

\subsection{A finite formula at every moment order}

Our double-zeta convention is
\[
 \zeta(a,b)=\sum_{\ell>k\ge1}\ell^{-a}k^{-b},
 \qquad a\ge2,\quad b\ge1.
\]
Only convergent double zeta values occur in the formulas below.
For a positive odd integer \(d\), define the ordered contribution
\begin{equation}\label{tail:ordered-block}
 U_{p,q}(d)=
 \begin{cases}
 \displaystyle
 \zeta(p,q-d)+\frac12\zeta(W-d),&d<q,\\[6pt]
 \displaystyle
 \sum_{h=1}^{\min\{(d-q+1)/2,\,p/2-1\}}
 \frac{\binom{d-q+1}{2h}}{d-q+1}
 B_{d-q+1-2h}\zeta(p-2h),&d>q.
 \end{cases}
\end{equation}
The two cases exhaust the possibilities since \(q\) is even.
The upper limit in the second line is an integer; the sum may be
empty. Define also the rational contribution
\begin{equation}\label{tail:rational-block}
 R_{p,q}(d)=
 \begin{cases}
 0,&d<W-1,\\[3pt]
 \displaystyle
 -\frac{B_{d-W+1}}2
 \left\{\frac1p\binom{d-q}{p-1}
          +\frac1q\binom{d-p}{q-1}\right\},&d\ge W-1,
 \end{cases}
\end{equation}
and
\begin{equation}\label{tail:block}
 C_{p,q}(d)=U_{p,q}(d)+U_{q,p}(d)+R_{p,q}(d).
\end{equation}
These expressions use only finite sums of explicitly specified
Bernoulli numbers and ordinary or double zeta values.

\begin{theorem}[All even orders and all centered moments]
\label{tail:all-moments}
For even \(p,q\ge2\) and every integer \(r\ge0\),
\begin{equation}\label{tail:moment-formula}
 \boxed{
 D_{p,q}(-2r)=
 \sum_{j=0}^{r}\frac{\binom{2r}{2j}}{2j+1}
 B_{2r-2j}(1/2)\,C_{p,q}(2j+1).}
\end{equation}
In particular this formula evaluates the ordinary sums in
\eqref{tail:ordinary-formula}.
\end{theorem}

\begin{proof}
Write \(\operatorname{CT}_N\) for the coefficient of
\(N^0(\log N)^0\) in an asymptotic expansion at positive infinity.
This coefficient will only be taken after exact finite summation.
Let
\begin{equation}\label{tail:S}
 S_r(k)=\sum_{n=0}^{k-1}(n+1/2)^{2r}
       =\frac{B_{2r+1}(k+1/2)}{2r+1}
       =\sum_{j=0}^{r}
         \frac{\binom{2r}{2j}}{2j+1}
         B_{2r-2j}(1/2)k^{2j+1}.
\end{equation}
This is an odd polynomial in \(k\), with zero constant coefficient.
Consequently every polynomial partial sum of a subtracted even power
in \eqref{tail:ordinary-formula} has zero constant coefficient in
\(N\). It follows that
\begin{equation}\label{tail:CT-moment}
 D_{p,q}(-2r)=\operatorname{CT}_N
              \sum_{n=0}^{N-1}x_n^{2r}A_p(n)A_q(n).
\end{equation}

The identity \(A_p(k-1)=A_p(k)+k^{-p}\) yields finite summation by
parts:
\begin{align}\label{tail:finite-summation}
 \sum_{n=0}^{N-1}x_n^{2r}A_p(n)A_q(n)
 ={}&S_r(N)A_p(N)A_q(N)\notag\\
 &+\sum_{k=1}^{N}S_r(k)
  \left\{\frac{A_p(k)}{k^q}+\frac{A_q(k)}{k^p}
                     +\frac1{k^{p+q}}\right\}.
\end{align}
For a positive odd \(d\), the coefficient block to evaluate is
\begin{equation}\label{tail:E}
 E_{p,q,d}(N)=N^dA_p(N)A_q(N)
 +\sum_{k=1}^{N}k^{d-q}A_p(k)
 +\sum_{k=1}^{N}k^{d-p}A_q(k)
 +\sum_{k=1}^{N}k^{d-W}.
\end{equation}
We prove that its constant coefficient is \(C_{p,q}(d)\).

First suppose \(d<q\). The corresponding tail sum converges, and
its limit is exactly
\begin{equation}\label{tail:negative-power}
 \sum_{k\ge1}k^{d-q}A_p(k)=\zeta(p,q-d).
\end{equation}
No continued or divergent double sum is used. If \(m=d-q\ge1\),
then \(m\) is odd. Put
\[
 F_m(N)=\sum_{k=1}^{N}k^m
       =\frac{B_{m+1}(N+1)-B_{m+1}}{m+1}.
\]
Interchanging a finite sum gives the exact identity
\begin{equation}\label{tail:Faulhaber-change}
 \sum_{k=1}^{N}k^mA_p(k)
 =A_p(N)F_m(N)
   +\sum_{\ell=1}^{N}\frac{F_m(\ell-1)}{\ell^p}.
\end{equation}
Since
\[
 F_m(\ell-1)=\frac1{m+1}
    \sum_{a=1}^{m+1}\binom{m+1}{a}B_{m+1-a}\ell^a,
\]
the constant term of the last sum is obtained directly from finite
power sums. Exponents \(a-p\ge0\) give zero constant; exponent
\(-1\) gives \(\gamma\); smaller exponents give \(\zeta(p-a)\).
Because \(m\) is odd, the nonzero Bernoulli coefficients have
\(a\) even, except for the half-leading term \(a=m\), whose
coefficient is \(-1/2\). Thus the even values contribute
\begin{equation}\label{tail:even-zeta-block}
 \sum_{h=1}^{\min\{(m+1)/2,p/2-1\}}
 \frac{\binom{m+1}{2h}}{m+1}B_{m+1-2h}\zeta(p-2h).
\end{equation}
If \(d<W-1\), the half-leading term contributes
\(-\zeta(W-d)/2\). At \(d=W-1\), it instead contributes
\(-\gamma/2\). If \(d>W-1\), it contributes zero.

The final power sum in \eqref{tail:E} contributes
\(\zeta(W-d)\) for \(d<W-1\), \(\gamma\) for \(d=W-1\), and
zero for \(d>W-1\). Combining it with the two half-leading terms
has a useful consequence. Each ordered block with \(d<q\) retains
one half of \(\zeta(W-d)\), while every block with \(d>q\)
retains none. At \(d=W-1\), the two terms \(-\gamma/2\) cancel
the final \(+\gamma\) exactly. Equations
\eqref{tail:negative-power} and \eqref{tail:even-zeta-block}
therefore give precisely \(U_{p,q}(d)+U_{q,p}(d)\).

It remains to compute the rational boundary contributions. The
uncentered expansion is
\begin{equation}\label{tail:uncentered}
 A_p(N)\sim\frac{N^{1-p}}{p-1}-\frac12N^{-p}
 +\sum_{j\ge1}\frac1{p-1}
       \binom{p+2j-2}{p-2}B_{2j}N^{1-p-2j}.
\end{equation}
For any odd integer \(k\), define the inverse-power coefficient by
\[
 u_p(k)=
 \begin{cases}
 0,&k<p-1,\\[2pt]
 \displaystyle\frac1{p-1}\binom{k-1}{p-2}B_{k-p+1},&k\ge p-1.
 \end{cases}
\]
The only nonzero even inverse-power coefficient in
\eqref{tail:uncentered} is \(-1/2\), at exponent \(p\).
In the Faulhaber polynomial \(F_m(N)\), all powers are even
except for its half-leading term \(N^m/2\). Pairing the possible
exponents therefore gives, for odd \(m\ge1\),
\begin{equation}\label{tail:single-boundary}
 \operatorname{CT}_N\{A_p(N)F_m(N)\}
       =\frac{p-m}{2p}\,u_p(m).
\end{equation}
Indeed the two potential contributions are
\(u_p(m)/2\) and
\(-\binom{m+1}{p}B_{m+1-p}/(2(m+1))\); the latter equals
\(-m u_p(m)/(2p)\). If \(m<p-1\), both are zero.

The odd coefficient of the product in the first term of
\eqref{tail:E} is
\[
 \operatorname{CT}_N\{N^dA_p(N)A_q(N)\}
      =-\frac12\{u_p(d-q)+u_q(d-p)\},
\]
where a term below its support is zero. Adding this to the two
contributions \eqref{tail:single-boundary} gives
\[
 -\frac{d-q}{2p}u_p(d-q)
 -\frac{d-p}{2q}u_q(d-p).
\]
It vanishes for \(d<W-1\). Otherwise, the binomial identity
\[
 \frac{d-q}{p-1}\binom{d-q-1}{p-2}
          =\binom{d-q}{p-1}
\]
reduces it to \(R_{p,q}(d)\) in \eqref{tail:rational-block}.
We have proved \(\operatorname{CT}_NE_{p,q,d}(N)=C_{p,q}(d)\).
Combining \eqref{tail:S}, \eqref{tail:CT-moment}, and
\eqref{tail:finite-summation} proves the theorem.
\end{proof}

\subsection{Reduction to a fixed set of ordinary zeta coordinates}

There is no unresolved depth-two constant in the preceding theorem.
Every double zeta appearing in \eqref{tail:ordered-block} has even
outer order, odd inner order, and odd total weight. The classical
odd-weight reduction proved by Borwein, Borwein, and Girgensohn
\cite{BBG} gives the following explicit elimination. For even
\(a\ge2\), odd \(b\ge3\), and \(w=a+b\),
\begin{align}\label{tail:Euler-reduction}
 \zeta(a,b)={}&\frac12\left\{\binom wa-1\right\}\zeta(w)
              +\zeta(a)\zeta(b)\notag\\
 &-\sum_{k=1}^{(w-3)/2}
 \left\{\binom{w-2k-1}{b-1}+\binom{w-2k-1}{a-1}\right\}
                 \zeta(2k)\zeta(w-2k).
\end{align}
For the boundary case \(b=1\), Euler's formula is
\begin{equation}\label{tail:Euler-boundary}
 \zeta(a,1)=\frac a2\zeta(a+1)
            -\sum_{k=1}^{a/2-1}\zeta(2k)\zeta(a+1-2k).
\end{equation}
All zeta arguments on the right are at least two. These are classical
inputs, not novelty claims of this article.

\begin{corollary}[A moment-independent coordinate space]
\label{tail:fixed-space}
Let \(M=\max\{p,q\}\). Every value \(D_{p,q}(-2r)\), for
arbitrary \(r\ge0\), lies in the same rational span
\begin{equation}\label{tail:fixed-space-formula}
 \operatorname{span}_{\mathbb Q}
 \left(\{1\}\cup
       \{\zeta(2h):1\le h\le M/2-1\}\cup
       \{C_{p,q}(d):1\le d<M,\ d\text{ odd}\}\right).
\end{equation}
This is a spanning set of \(M\) specified numbers, each explicitly
reducible to ordinary zeta values and their pairwise products.
It is not a claim that those numbers are linearly independent.
\end{corollary}

\begin{proof}
For every odd \(d>M\), both ordered blocks in
\eqref{tail:ordered-block} contain only even zeta values of argument
at most \(M-2\), and the remaining block is rational. The odd
integers \(d<M\) give the finite list in
\eqref{tail:fixed-space-formula}. The moment formula proves the
inclusion, while \eqref{tail:Euler-reduction} and
\eqref{tail:Euler-boundary} give the stated ordinary-zeta reduction.
\end{proof}

This finiteness is uniform in the moment order. Increasing \(r\)
changes only rational Bernoulli coefficients; it introduces no new
arithmetic constants.

\subsection{The first mixed pair: a complete \texorpdfstring{$(2,4)$}{(2,4)} evaluation}

For \((p,q)=(2,4)\), write \(D_r=D_{2,4}(-2r)\). In this case
the entire block sequence has the short form
\begin{align}\label{tail:24-blocks}
 C_{2,4}(1)&=\frac{15}{2}\zeta(5)-3\zeta(2)\zeta(3),\notag\\
 C_{2,4}(3)&=\frac32\zeta(3)+\frac12\zeta(2),\notag\\
 C_{2,4}(2j+1)&=\frac{2j-1}{2}B_{2j-2}\zeta(2)
  -\frac{(2j-3)(2j^2-3j+7)}{24}B_{2j-4},\qquad j\ge2.
\end{align}
For example, the first line follows from
\(\zeta(2,3)=9\zeta(5)/2-2\zeta(2)\zeta(3)\) and
\(\zeta(4,1)=2\zeta(5)-\zeta(2)\zeta(3)\). Thus the arithmetic
space in this case can be written especially economically as
\begin{equation}\label{tail:24-space}
 \operatorname{span}_{\mathbb Q}
 \{1,\zeta(2),\zeta(3),5\zeta(5)-2\zeta(2)\zeta(3)\}.
\end{equation}

The first four values are
\begin{align}\label{tail:24-values}
 D_0={}&\frac{15}{2}\zeta(5)-3\zeta(2)\zeta(3),\notag\\
 D_1={}&-\frac58\zeta(5)+\frac14\zeta(2)\zeta(3)
                    +\frac12\zeta(3)+\frac16\zeta(2),\notag\\
 D_2={}&\frac7{32}\zeta(5)-\frac7{80}\zeta(2)\zeta(3)
              -\frac14\zeta(3)-\frac1{30}\zeta(2)-\frac3{40},\notag\\
 D_3={}&-\frac{155}{896}\zeta(5)+\frac{31}{448}\zeta(2)\zeta(3)
              +\frac7{32}\zeta(3)-\frac1{672}\zeta(2)+\frac{31}{672}.
\end{align}
The first two sums converge without subtraction. The next two yield
the polygamma identities
\begin{align}\label{tail:24-polygamma-sums}
 \sum_{n=0}^{\infty}
 \left\{\frac{x_n^4}{6}\psi'(n+1)\psi'''(n+1)-\frac13\right\}
       &=D_2,\notag\\
 \sum_{n=0}^{\infty}
 \left\{\frac{x_n^6}{6}\psi'(n+1)\psi'''(n+1)
                   -\frac{x_n^2}{3}+\frac7{36}\right\}
       &=D_3.
\end{align}
The signs and normalizations of both subtractions follow directly
from \(k_{2,4,0}=1/3\) and \(k_{2,4,1}=-7/36\).

There is also an exact finite-generation statement. The inverse
Bernoulli transform is
\begin{equation}\label{tail:inverse-transform}
 C_{p,q}(2j+1)=\sum_{r=0}^{j}
          \binom{2j+1}{2r}2^{-2(j-r)}D_{p,q}(-2r).
\end{equation}
It follows by multiplying the generating function in the next
section by \(2\sinh(t/2)/t\), or directly by the Bernoulli
generating identity. Put \(C_j=C_{2,4}(2j+1)\). Then
\[
 C_2=\frac14\zeta(2)-\frac38,
 \qquad C_3=-\frac1{12}\zeta(2)-\frac13,
\]
so
\[
 1=-\frac8{11}(C_2+3C_3),\qquad
 \zeta(2)=\frac{32C_2-36C_3}{11},\qquad
 \zeta(3)=\frac23C_1-\frac{32}{33}C_2+\frac{12}{11}C_3.
\]
Consequently \emph{every} \(D_r\) is an explicit rational linear
combination of \(D_0,D_1,D_2,D_3\). This is an exact generation
statement, independent of any conjectured arithmetic independence.
Two useful eliminations are
\begin{align}\label{tail:24-eliminations}
 7D_0+412D_1+2000D_2+1344D_3&=-88,\notag\\
 23D_0+524D_1-2480D_2-4032D_3&=176\zeta(2).
\end{align}
For instance the first is the single absolutely convergent identity
\begin{equation}\label{tail:rational-polygamma}
 \sum_{n=0}^{\infty}
 \left\{
 (1344x_n^6+2000x_n^4+412x_n^2+7)
       \frac{\psi'(n+1)\psi'''(n+1)}6
       -448x_n^2-\frac{1216}{3}
 \right\}=-88.
\end{equation}


% END sections/mixed_even_tails.tex


% BEGIN sections/mixed_even_tail_generators.tex
\section{A classical-polylogarithm generator for every pair}
\label{tail:generator-section}

The finite formula of Theorem~\ref{tail:all-moments} can be summed
at every moment order. This yields an analytic identity, as opposed
to a merely formal generating series.

For integers \(a\ge2\) and \(b\ge1\), define
\begin{equation}\label{tail:I}
 I_{a,b}(t)=\int_0^t u^a(t-u)^b\coth(u/2)\,\dd u.
\end{equation}
Inside \(|t|<2\pi\), use the straight segment from zero to \(t\).
The integrand has a removable zero-endpoint singularity, and the
integral is unambiguous there. Define the finite polynomial
\begin{equation}\label{tail:P}
 P_{p,q}(t)=\sum_{\substack{1\le d<q\\d\ {
 \rm odd}}}
 \left\{\zeta(p,q-d)+\frac12\zeta(W-d)\right\}
                        \frac{t^{d-1}}{d!}.
\end{equation}
Every double zeta in this polynomial is eliminated by
\eqref{tail:Euler-reduction} or \eqref{tail:Euler-boundary}.

\begin{theorem}[All-pair analytic moment generator]
\label{tail:generator}
For even \(p,q\ge2\), the series
\[
 G_{p,q}(t)=\sum_{r=0}^{\infty}
           D_{p,q}(-2r)\frac{t^{2r}}{(2r)!}
\]
converges in \(|t|<2\pi\). In that disk it is the even holomorphic
function
\begin{equation}\label{tail:generator-formula}
 \boxed{
 G_{p,q}(t)=\frac{t}{2\sinh(t/2)}
             \{F_{p,q}(t)+F_{q,p}(t)\},}
\end{equation}
where all singularities at zero are removable and
\begin{align}\label{tail:F}
 F_{p,q}(t)={}&P_{p,q}(t)\notag\\
 &+\sum_{h=1}^{p/2-1}
   \frac{\zeta(p-2h)}{2t(2h)!(q-1)!}
       I_{2h,q-1}(t)
  -\frac{I_{p,q-1}(t)}{4t p!(q-1)!}.
\end{align}
For real \(0<t<2\pi\), the expression in braces has a finite
form using only rational functions of \(t\), ordinary zeta values,
and \(\operatorname{Li}_k(e^{-t})\) for \(2\le k\le p+q\).
The full generator includes the displayed elementary prefactor
\(t/(2\sinh(t/2))\).
The integral formula fixes its holomorphic continuation at zero.
\end{theorem}

\begin{proof}
The Bernoulli convolution in \eqref{tail:moment-formula} is
equivalent formally to
\begin{equation}\label{tail:block-generator}
 G_{p,q}(t)=\frac{t}{2\sinh(t/2)}
           \sum_{\substack{d\ge1\\d\ {\rm odd}}}
                        C_{p,q}(d)\frac{t^{d-1}}{d!}.
\end{equation}
We show that the last series converges in the stated disk and
equals \(F_{p,q}+F_{q,p}\).

Write
\[
 E(u)=\frac u2\coth(u/2)
     =\sum_{k\ge0}\frac{B_{2k}}{(2k)!}u^{2k},
 \qquad |u|<2\pi.
\]
The terms with \(d<q\) in the ordered block are exactly
\(P_{p,q}\). For each fixed \(1\le h\le p/2-1\), put
\(d=q+2h-1+2k\). Its remaining Bernoulli contribution is
\[
 \frac{\zeta(p-2h)}{(2h)!}
 \sum_{k\ge0}\frac{(2h+2k-1)!}{(q+2h+2k-1)!}
            \frac{B_{2k}}{(2k)!}t^{q+2h+2k-2}.
\]
Euler's beta integral transforms this normally convergent series
into
\[
 \frac{\zeta(p-2h)}{t(2h)!(q-1)!}
       \int_0^t u^{2h-1}(t-u)^{q-1}E(u)\,\dd u
   =\frac{\zeta(p-2h)I_{2h,q-1}(t)}
          {2t(2h)!(q-1)!}.
\]
For the first binomial term in \(R_{p,q}(d)\), set
\(d=W-1+2k\). Its series is
\[
 -\frac1{2p!}\sum_{k\ge0}
    \frac{(p+2k-1)!}{(W+2k-1)!}
         \frac{B_{2k}}{(2k)!}t^{W+2k-2},
\]
which the same beta integral turns into
\(-I_{p,q-1}(t)/(4t p!(q-1)!)\). The second term is the
same expression with \(p,q\) exchanged. This proves
\eqref{tail:F} and \eqref{tail:generator-formula}.

All power-series integrations above are normal on compact subsets
of \(|t|<2\pi\). The prefactor \(t/(2\sinh(t/2))\) is
holomorphic there, including at zero. The resulting function is
even: this follows from its Taylor construction, or by substituting
\(u=tv\) in each integral. Thus
\eqref{tail:block-generator} is the Taylor series of an even
holomorphic function, proving the claimed convergence.

Finally, write \(\coth(u/2)=1+2/(e^u-1)\) and expand the
polynomial \((t-u)^b\). For \(t>0\), elementary integration
gives
\begin{align}\label{tail:finite-Li}
 I_{a,b}(t)={}&\frac{a!b!}{(a+b+1)!}t^{a+b+1}\notag\\
 &+2\sum_{\ell=0}^{b}(-1)^\ell\binom b\ell
                     t^{b-\ell}J_{a+\ell}(t),\notag\\
 J_m(t)={}&m!\left\{\zeta(m+1)
       -\sum_{j=0}^{m}\frac{t^j}{j!}
                  \operatorname{Li}_{m+1-j}(e^{-t})\right\}.
\end{align}
One proof expands \((e^u-1)^{-1}\) into its positive exponential
series and integrates each term. Since \(m\ge2\), the
integrals are absolutely convergent. The coefficient of
\(\operatorname{Li}_1(e^{-t})\) in the first two lines is
proportional to \(\sum_{\ell=0}^{b}(-1)^\ell\binom b\ell=0\).
The largest remaining polylogarithm order is
\(a+b+1\le p+q\). This proves the last assertion.
\end{proof}

\subsection{An explicit \texorpdfstring{$\operatorname{Li}_2$--$\operatorname{Li}_6$}{Li2--Li6} identity}

For \((p,q)=(2,4)\), the general generator becomes
\begin{align}\label{tail:24-integral-generator}
 G_{2,4}(t)=\frac{t}{2\sinh(t/2)}\bigg\{
 &\frac{15}{2}\zeta(5)-3\zeta(2)\zeta(3)
       +\frac{t^2}{4}\zeta(3)
       +\frac{\zeta(2)}{4t}I_{2,1}(t)\notag\\
 &-\frac{I_{4,1}(t)}{96t}
       -\frac{I_{2,3}(t)}{48t}\bigg\}.
\end{align}
Put \(L_k(t)=\operatorname{Li}_k(e^{-t})\). An entirely explicit
version, valid for real \(0<t<2\pi\), is
\begin{equation}\label{tail:24-Li-generator}
 G_{2,4}(t)=\frac{t}{2\sinh(t/2)}\,\mathcal L_{2,4}(t),
\end{equation}
where
\begin{align}\label{tail:24-Li-body}
 \mathcal L_{2,4}(t)={}&4\zeta(5)-2\zeta(2)\zeta(3)
       +\frac{t^2}{6}\zeta(3)+\frac{3t}{4}\zeta(4)
       +\frac{t^3}{48}\zeta(2)-\frac{t^5}{1440}\notag\\
 &+\left(\frac{t}{2}\zeta(2)-\frac{t^3}{48}\right)L_2(t)
       +\left(2\zeta(2)-\frac{t^2}{6}\right)L_3(t)\notag\\
 &-tL_4(t)-4L_5(t)
       +\frac{3\zeta(2)}{t}\{L_4(t)-\zeta(4)\}
       +\frac{15}{2t}\{\zeta(6)-L_6(t)\}.
\end{align}
Although individual polylogarithms have logarithmic expansions at
zero, their specified combination is even and holomorphic there.
Equation~\eqref{tail:24-integral-generator}, rather than a separate
choice of logarithm for each term, provides the continuation. The
coefficient of \(t^{2r}/(2r)!\) is exactly the all-order moment in
\eqref{tail:moment-formula} and \eqref{tail:24-blocks}.

\subsection{Scope and next exact questions}

The two-tail problem is closed by a finite Bernoulli formula,
ordinary-zeta elimination, and a convergent classical-polylogarithm
generator. The next questions are structurally different.

\begin{enumerate}
\item For four even tails, determine the smallest useful
multiple-zeta coordinate space uniform in the centered moment
order. Finite summation by parts still lowers the number of tails,
but higher-depth constants can enter. A two-tail calculation does
not prove their elimination.
\item Determine exact first spectral derivatives at the regular
points \(s=-2r\). These contain logarithmically weighted remainder
sums in addition to local Bernoulli data. The value theorem alone
does not reduce those global contributions.
\item Introduce rational inner shifts while retaining a matched
center for the outer weight. Identify the resulting finite
colored Euler-sum coordinates and the precise analogues of the
classical-polylogarithm generator.
\item For fixed \(p,q\), find optimal finite generating sets of
the actual moment sequence. Corollary~\ref{tail:fixed-space}
provides at most \(\max\{p,q\}\) explicit coordinates, and
\((2,4)\) has an exact four-moment generator. Arithmetic
independence of a proposed minimal set must be distinguished from
the rank of its rational coefficient matrix.
\end{enumerate}


% END sections/mixed_even_tail_generators.tex


% BEGIN sections/complex_powers.tex
\section{Complex powers of a cotangent resolvent}
\label{power:sec}

The rational resolvent theorem in the incoming harmonic-parity report
leaves open its extension to complex powers (Question R4).  The
complex-power sine transform in the higher-reflection report treats a
different, reflection-symmetric weight.  We now give a one-sided
complex-power transform at every nonreal resolvent parameter, an ordinary
Euler integral involving a single classical polylogarithm, and a completed
germ at every integral power.  The latter includes a removable but
nonzero endpoint correction when the integral power is negative.
Classical Hurwitz Fourier analysis and beta integrals supply the analytic
tools \cite{DLMF,EspinosaMoll}; the formulas and their
normalizations below are proved directly.

\subsection{A branch fixed by the interval}

For $z\in\mathbb C\setminus\mathbb R$, put
\begin{equation}\label{power:coordinates}
 \epsilon=\operatorname{sgn}\Im z,\quad
 a=z+\epsilon\pi i,\quad q=\frac{z-\epsilon\pi i}{a},\quad
 \sigma=-2\epsilon\pi i,\quad C=-\frac1a=\frac{1-q}{\sigma}.
\end{equation}
Thus $|q|<1$.  Throughout,
\[
 \Log\sigma=\log(2\pi)-\epsilon\pi i/2,
 \qquad \Log C=\Log(1-q)-\Log\sigma.
\]
Both $\Log C$ and $\Log(1-q)$ in these equations are principal
logarithms.  Write
\begin{equation}\label{power:R}
 K(x)=\pi\cot(\pi x),\quad R_z(x)=\frac1{K(x)-z},\quad
 w=e^{2\epsilon\pi ix}.
\end{equation}
Direct algebra gives
\begin{equation}\label{power:factor}
 R_z(x)=C\frac{1-w}{1-qw}.
\end{equation}
The continuous logarithm of $R_z(x)$ on $0<x<1$ is fixed by
$\Log R_z(x)-\log x\to0$ at zero.  Its imaginary part lies between
$0$ and $\epsilon\pi$, and it approaches $\epsilon\pi$ at one.
In particular,
\begin{equation}\label{power:branches}
 R_z(x)^\lambda=\exp(\lambda\Log R_z(x)),\qquad
 R_z(1-y)^\lambda=e^{\epsilon\pi i\lambda}y^\lambda
                       V_z(-y)^\lambda,
 \quad V_z(x)=\frac{R_z(x)}x.
\end{equation}
Here $V_z$ is analytic near zero and $V_z(0)=1$.  Powers of $V_z$
use the analytic logarithm vanishing at zero.  These prescriptions also
specify all logarithmic derivatives in $\lambda$.

\subsection{A beta mixture of ordinary polylogarithms}

Define coefficients and their Dirichlet transform by
\begin{align}
 \left(\frac{1-w}{1-qw}\right)^\lambda
   &=\sum_{n\ge0}d_n(\lambda,q)w^n,\quad d_0=1,
       \label{power:dn}\\
 d_n(\lambda,q)&=\sum_{j=0}^n
       \frac{(-\lambda)_j(\lambda)_{n-j}}{j!(n-j)!}q^{n-j},
       \label{power:dnfinite}\\
 \mathcal L_\lambda(\nu,q)&=\sum_{n\ge1}
                               \frac{d_n(\lambda,q)}{n^\nu}.
       \label{power:L}
\end{align}
The last series is initially normally convergent for
$\Re(\nu+\lambda)>0$, on compact parameter sets away from the boundary
of $|q|<1$.  At integral $\lambda$ it can have a larger convergence
domain; our continuation will not rely on such exceptional improvements.
The notation in \eqref{power:L} is a transform of the explicit
coefficients \eqref{power:dnfinite}, not a claim that an unevaluated
integral has been reduced to familiar constants.

\begin{theorem}[The Euler--polylogarithm formula]\label{power:eulerthm}
For $-1<\Re\lambda<1$, $|q|<1$, and initially $\Re\nu>1$,
\begin{equation}\label{power:euler}
 \boxed{\displaystyle
 \mathcal L_\lambda(\nu,q)=
 -\frac{\sin\pi\lambda}{\pi}(1-q)
 \int_0^1 y^{-\lambda}(1-y)^\lambda
       \frac{\Li_\nu(q+(1-q)y)}{q+(1-q)y}\,\dd y.}
\end{equation}
The quotient $\Li_\nu(t)/t$ has value $1$ at $t=0$.  Thus the
formula remains valid when the segment from $q$ to $1$ passes through
zero.  Within the same $\lambda$ strip the formula also holds on
$\Re(\nu+\lambda)>0$ by analytic continuation of the convergent
integrals.  All fixed parameter derivatives commute with these integrals
on compact subsets of their convergence domains.
\end{theorem}

\begin{proof}
We first prove the coefficient identity
\begin{equation}\label{power:beta-coeff}
 d_n(\lambda,q)=-\frac{\sin\pi\lambda}{\pi}(1-q)
 \int_0^1y^{-\lambda}(1-y)^\lambda
                 (q+(1-q)y)^{n-1}\,\dd y,\qquad n\ge1.
\end{equation}
For $0<q<1$, deform the Cauchy coefficient contour for
$((1-w)/(1-qw))^\lambda$ onto $1<w<1/q$.  The difference of its
upper and lower boundary values is
$-2i\sin(\pi\lambda)((w-1)/(1-qw))^\lambda$.
The endpoint integrals exist precisely in the stated strip.  The
integral at infinity vanishes for $n\ge1$.  Substituting $t=1/w$
and then $t=q+(1-q)y$ proves \eqref{power:beta-coeff}.
Alternatively, its right side can be expanded as a finite beta sum
and checked against \eqref{power:dnfinite}.  Both sides are
polynomials in $q$, so the identity holds for all complex $q$.

For $\Re\nu>1$ we may sum \eqref{power:beta-coeff} after multiplying
by $n^{-\nu}$.  The segment lies in the closed unit disk and only
its last endpoint can meet the unit circle.  The absolutely convergent
polylogarithm series, together with the integrable beta weight, provides
domination.  This yields \eqref{power:euler}.  Near $y=1$ the possible
singular part of $\Li_\nu(q+(1-q)y)$ has order
$(1-y)^{\nu-1}$, with the usual logarithmic limits at integral
orders \cite[Eq.~25.12.12]{DLMF}.  The conditions
$\Re\lambda>-1$ and $\Re(\nu+\lambda)>0$ therefore give the
asserted extension and the same domination after any fixed number of
derivatives.  Near $y=0$, $\Re\lambda<1$ suffices, including $q=0$.
\end{proof}

There is a useful independent Mellin representation:
\begin{equation}\label{power:mellinL}
 \mathcal L_\lambda(\nu,q)=\frac1{\Gamma(\nu)}
 \int_0^\infty t^{\nu-1}
 \left\{\left(\frac{1-e^{-t}}{1-qe^{-t}}\right)^\lambda-1\right\}
 \,\dd t,
\end{equation}
initially for $\Re\nu>\max(0,-\Re\lambda)$.
At zero, the expression in braces is $t^\lambda$ times an analytic
function minus $1$; at infinity it decreases exponentially.
Finite Taylor subtraction in \eqref{power:mellinL} proves joint
meromorphic continuation in $(\nu,\lambda)$.  Its only possible
moving polar hyperplanes are
$\nu+\lambda=-j$ for $j\in\mathbb Z_{\ge0}$; some disappear
through zeros of $1/\Gamma(\nu)$ or of their coefficients.
This supplies a continuation independently of the Euler integral's
restricted $\lambda$ strip.

\subsection{The Hurwitz transform and an ordinary log-Gamma integral}

\begin{theorem}[Complex-power Hurwitz transform]\label{power:hurwitzthm}
For $p\in\mathbb Z_{\ge0}$, $\Re\lambda>-1$, and
$\Re(u+\lambda)>p$, away from the Hurwitz pole when $p=0$,
\begin{equation}\label{power:hurwitz}
 \boxed{\displaystyle
 \int_0^1\partial_x^p\zeta(1-u,x)R_z(x)^\lambda\,\dd x
 =C^\lambda\Gamma(u)\sigma^{p-u}
                    \mathcal L_\lambda(u-p,q).}
\end{equation}
The identity then continues meromorphically in its parameters.  Also,
for every $\Re\lambda>-1$,
\begin{equation}\label{power:mean}
 \int_0^1R_z(x)^\lambda\,\dd x=C^\lambda.
\end{equation}
\end{theorem}

\begin{proof}
For $\Re u>p+1$, the Hurwitz Fourier series and its first $p$
argument derivatives are absolutely convergent.  The nonzero Fourier
pairings are
\[
 \int_0^1\partial_x^p\zeta(1-u,x)e^{2\epsilon\pi inx}\,\dd x
     =\Gamma(u)\sigma^{p-u}n^{p-u},\qquad n\ge1,
\]
and the zero mode vanishes.  Multiply by the Abel-regularized expansion
\eqref{power:dn} and then let the Abel radius tend to one.
For $\Re\lambda>-1$ its boundary values are uniformly integrable:
near either endpoint they are bounded by a constant times
$x^{\min(\Re\lambda,0)}$ or $(1-x)^{\min(\Re\lambda,0)}$.
This proves the formula in an initial open region.  The Dirichlet
series is normally convergent under $\Re(u+\lambda)>p$.
The original integral has the same sufficient local conditions, since
its only nonsmooth Hurwitz term at zero is
$\prod_{j=1}^p(u-j)x^{u-p-1}$.  Holomorphic continuation within
that domain proves \eqref{power:hurwitz}.  The constant Fourier
coefficient of the Abel kernel is $C^\lambda$; the same integrable
limit proves \eqref{power:mean}.
\end{proof}

\begin{corollary}[An ordinary Gamma moment]\label{power:gamma}
For $\Re\lambda>-1$,
\begin{align}\label{power:gammaformula}
 \int_0^1\log\Gamma(x)R_z(x)^\lambda\,\dd x
 =C^\lambda\left\{\frac{\log(2\pi)}2
 +\frac{(\gamma+\Log\sigma)\mathcal L_\lambda(1,q)
       -\partial_\nu\mathcal L_\lambda(1,q)}{\sigma}\right\}.
\end{align}
\end{corollary}
\begin{proof}
Differentiate \eqref{power:hurwitz} with $p=0$ at $u=1$ and use
$\zeta'(0,x)=\log\Gamma(x)-\tfrac12\log(2\pi)$ and
$\Gamma'(1)=-\gamma$ \cite{DLMF}.
The logarithmic endpoint singularity is integrable under the stated
condition, so this is an ordinary integral and an ordinary derivative.
\end{proof}

Thus, in the strip of Theorem~\ref{power:eulerthm}, the Gamma moment
is an explicit combination of two integrals of $\Li_1$ and its first
order derivative against an elementary beta weight.  At $q=0$, its
first transform is $\mathcal L_\lambda(1,0)=-\psi(1+\lambda)-\gamma$.
Section~\ref{moment:sec} gives the stronger finite Gamma generating
formula for all polynomial moments.

\subsection{Completion at every integral power}

Let
\begin{equation}\label{power:stieltjes}
 g(u,x)=\zeta(1-u,x)+\frac1u
     =\sum_{m\ge0}\gamma_m(x)\frac{u^m}{m!},\qquad
 A_p(u)=\prod_{j=1}^p(u-j),\quad A_0=1.
\end{equation}
Define finite local coefficients
\begin{align}
 v_j(\lambda;z)&=[x^j]V_z(x)^\lambda,\label{power:vj}\\
 h_p(u,x)&=\partial_x^p\left\{\zeta(1-u,1+x)+\frac1u\right\},
        \label{power:hp}\\
 b_j(u,\lambda;z)&=[x^j]h_p(u,x)V_z(x)^\lambda.
        \label{power:bj}
\end{align}
The function $h_p$ is holomorphic near $(u,x)=(0,0)$.  Its coefficients
are explicit: the coefficient of $x^j$ is
\begin{equation}\label{power:hcoeff}
 \frac{(-1)^{p+j}(1-u)_{p+j}}{j!}\zeta(1-u+p+j)
 \quad(p+j>0),
\end{equation}
and is $\zeta(1-u)+1/u$ when $p=j=0$.
The first coefficients of $V_z$ are
\[
 V_z(x)=1+zx+(z^2+\pi^2/3)x^2
                   +(z^3+2\pi^2z/3)x^3+O(x^4).
\]
All fixed Taylor jets of $v_j$ and $b_j$ required for coefficient
extraction are therefore finite polynomials in $z$, the expansion
parameters, zeta values or their derivatives, and Stieltjes
coefficients.  No new global integral enters them.

For an integer $L$, set $\lambda=L+v$ and
\begin{align}
 \mathcal T_p(u,\lambda;z)
   &=C^\lambda\Gamma(u)\sigma^{p-u}
                    \mathcal L_\lambda(u-p,q)
         +\mathbf1_{p=0}\frac{C^\lambda}{u},\label{power:T}\\
 \mathcal H_{p,L}(u,v;z)
   &=\mathcal T_p(u,L+v;z)
      -\mathbf1_{L\le p}
        \frac{A_p(u)v_{p-L}(L+v;z)}{u+v}\notag\\
   &\quad-\mathbf1_{L\le-1}
       \frac{\{1+(-1)^{-L-1}e^{\epsilon\pi i(L+v)}\}
                    b_{-L-1}(u,L+v;z)}{v}.
       \label{power:H}
\end{align}
A term with a false indicator is omitted completely.  In the last
line the quotient is read by its removable analytic extension.

\begin{theorem}[Joint Stieltjes and power completion]\label{power:completion}
For each $p\ge0$, $L\in\mathbb Z$, and $z\notin\mathbb R$,
$\mathcal H_{p,L}$ is holomorphic near $(u,v)=(0,0)$.
For all integers $m,k\ge0$,
\begin{equation}\label{power:allmoments}
 \boxed{\displaystyle
 \FP\int_0^1\gamma_m^{(p)}(x)R_z(x)^L
                    (\Log R_z(x))^k\,\dd x
       =m!k![u^mv^k]\mathcal H_{p,L}(u,v;z).}
\end{equation}
The finite part is taken in the fixed coordinates $x$ and $1-x$.
All coefficient extractions and all $z$ derivatives are locally
holomorphic on each open half-plane.
\end{theorem}

\begin{proof}
The exact local splitting is
\[
 \partial_x^p g(u,x)=A_p(u)x^{u-p-1}+h_p(u,x).
\]
At zero, multiplication by $R_z(x)^\lambda=x^\lambda V_z(x)^\lambda$
produces the two families
\[
 A_p(u)v_j(\lambda;z)x^{u+\lambda-p-1+j},\qquad
 b_j(u,\lambda;z)x^{\lambda+j}.
\]
After setting $\lambda=L+v$, their integration denominators pass
through zero only for $j=p-L\ge0$ and $j=-L-1\ge0$, respectively.
They give the second line of \eqref{power:H} and the first endpoint's
contribution to its third line.

At the other endpoint put $y=1-x$.  The smooth Hurwitz factor is
exactly $h_p(u,-y)$ and \eqref{power:branches} applies.  Its
coefficient of $y^j$ is
$e^{\epsilon\pi i\lambda}(-1)^jb_j(u,\lambda;z)$.
This supplies precisely the remaining term in the braces in
\eqref{power:H}.  Their sum vanishes on $v=0$, since
$(-1)^{-L-1}e^{\epsilon\pi iL}=-1$.  Nevertheless its quotient by
$v$ must be subtracted with the complete numerator.

To make the argument analytic, choose a small endpoint interval and
subtract finitely many of the displayed monomials, enough to leave
a uniformly integrable remainder throughout a small parameter
polydisk.  The middle interval is holomorphic.  Nonresonant local
denominators are bounded away from zero.  For a resonant monomial
with exponent $\alpha-1$, the difference between its meromorphic
integral and the generating function of its coordinate finite parts
is its full amplitude divided by $\alpha$, because
\[
 \frac{\delta^\alpha}{\alpha}
 =\frac1\alpha+\frac{\delta^\alpha-1}{\alpha}.
\]
The second summand is entire in $\alpha$ and its derivatives give
the constants of the corresponding logarithmic integrals.  This
also shows that the arbitrary splitting point cancels out.  Finally,
the global Hurwitz pole contributes $-C^\lambda/u$ only for $p=0$,
and \eqref{power:T} has already removed it by \eqref{power:mean}.
The remaining germ is holomorphic and its coefficients are exactly
\eqref{power:allmoments}.  The same normally convergent construction
is uniform on compact subsets of either $z$ half-plane.
\end{proof}

Away from integral powers there is no resonant exponent at $u=0$.
Consequently, for $\lambda\notin\mathbb Z$, the same local subtraction
proof gives the simpler identity
\begin{equation}\label{power:noninteger}
 \FP\int_0^1\gamma_m^{(p)}(x)R_z(x)^\lambda
             (\Log R_z(x))^k\,\dd x
 =m![u^m]\partial_\lambda^k\mathcal T_p(u,\lambda;z).
\end{equation}
Here the finite part means the constant term in the generalized
power--logarithm endpoint expansion.  For nonintegral $\lambda$,
none of its subtracted endpoint powers has exponent zero after
integration, so differentiation and this finite part commute locally
in $\lambda$.  At integral powers the full completion
\eqref{power:H}, rather than a limit of \eqref{power:noninteger},
specifies the coordinate finite part.

\paragraph{The correction at negative powers is not zero.}
Although the last quotient in \eqref{power:H} is removable, its
value at $v=0$ is
\begin{equation}\label{power:smoothcontact}
 -\epsilon\pi i\,b_{-L-1}(u,L;z),\qquad L\le-1.
\end{equation}
Removing only pole residues would miss it.  It records the different
endpoint phases of $R_z^\lambda$ before $\lambda$ is specialized to
an integer.  At a fixed negative integer the integrand itself is a
polynomial in $K(x)-z$, so the resulting finite parts must have the
corresponding polynomial dependence on $z$.  Formula
\eqref{power:smoothcontact} is necessary for that compatibility.

The same issue already appears without a Hurwitz factor.  For $M\ge1$,
\begin{equation}\label{power:negative-mean}
 \FP\int_0^1(K(x)-z)^M\,\dd x
 =\frac{(-z+i\pi)^M+(-z-i\pi)^M}{2},
\end{equation}
whereas the continued mean in \eqref{power:mean} is
$C^{-M}=(-z-\epsilon i\pi)^M$.  Reflection gives zero for odd
cotangent powers, and integration of their derivative recurrence gives
$\FP\int_0^1K^{2j}\,\dd x=(-\pi^2)^j$, proving
\eqref{power:negative-mean}.  The difference is exactly
$\epsilon\pi i\,[x^{M-1}]V_z(x)^{-M}$, as the two-endpoint
completion predicts.  This is an elementary check at every negative
integral power.
An independent all-$M$ residue proof is given in
Appendix~\ref{contact:verification}.

\begin{example}[Two normalization checks]\label{power:canaries}
For $L=0$ and $k=0$, the theorem gives
$\FP\int_0^1\gamma_m^{(p)}(x)\,\dd x=0$.
For $L=-1$, $p=m=k=0$, it gives
\begin{equation}\label{power:cotcheck}
 \FP\int_0^1\gamma_0(x)(\pi\cot\pi x-z)\,\dd x
        =-\frac{\pi^2}{2}.
\end{equation}
Indeed $\mathcal L_{-1}(u,q)=(1-q)\zeta(u)$,
$v_1(-1;z)=-z$, and $b_0=\zeta(1-u)+1/u$.
At $v=0$, \eqref{power:H} reduces, by the zeta functional equation, to
\[
 2\pi\Gamma(u)(2\pi)^{-u}\sin(\pi u/2)\zeta(u),
\]
whose value at zero is $-\pi^2/2$.  The same value follows directly
from digamma reflection and $\FP\int_0^1K(x)^2\,\dd x=-\pi^2$.
For $L=1$ and $p=1$, the moving term is $(u-1)/(u+v)$; at
$z=i\pi$ it recovers
\[
 \FP\int_0^1\frac{\psi'(x)}{\pi\cot\pi x-i\pi}\,\dd x
       =1-\gamma-\log(2\pi)+i\pi/2,
\]
including the inherited resolvent theorem's harmonic correction.
\end{example}

\subsection{An explicit logarithmic resolvent identity}

At first order in $\lambda$ there is no need to retain the transform
$\mathcal L$: differentiating \eqref{power:dn} at zero gives
\begin{equation}\label{power:firstlambda}
 \left.\partial_\lambda\mathcal L_\lambda(u,q)\right|_{\lambda=0}
       =\Li_{1+u}(q)-\zeta(1+u).
\end{equation}

\begin{corollary}[All Stieltjes indices with one logarithm]
\label{power:logcor}
The expression in braces below is holomorphic at $u=0$, and
\begin{equation}\label{power:alllog}
 \boxed{\displaystyle
 \FP\int_0^1\gamma_m(x)\Log R_z(x)\,\dd x
 =m![u^m]\left\{
 \Gamma(u)\sigma^{-u}\bigl(\Li_{1+u}(q)-\zeta(1+u)\bigr)
       +\frac{\Log C}{u}+\frac1{u^2}\right\}.}
\end{equation}
In particular, if a superscript $[1]$ denotes differentiation of the
polylogarithm with respect to its order, then
\begin{align}\label{power:logexplicit}
 \FP\int_0^1\gamma_0(x)\Log R_z(x)\,\dd x
  ={}&\gamma_1+\frac{\gamma^2}{2}-\frac{\zeta(2)}2
              -\frac{(\Log\sigma)^2}{2}\notag\\
    &+\Li_1^{[1]}(q)
               +(\gamma+\Log\sigma)\Log(1-q).
\end{align}
\end{corollary}
\begin{proof}
Differentiate \eqref{power:H} at $v=0$ with $p=L=0$, using
\eqref{power:firstlambda}.  The derivative of $-1/(u+v)$ is
$1/u^2$.  Expansion of
\[
 \Gamma(u)\sigma^{-u}
 =\frac1u\exp\left(-Au+\sum_{j\ge2}\frac{(-1)^j\zeta(j)}j u^j\right),
 \qquad A=\gamma+\Log\sigma,
\]
and $\zeta(1+u)=u^{-1}+\gamma-\gamma_1u+O(u^2)$ proves
\eqref{power:logexplicit}.  The coefficient of $u^{-1}$ cancels
because $\Li_1(q)=-\Log(1-q)$ and
$\Log C=\Log(1-q)-\Log\sigma$.
\end{proof}

At $z=\epsilon i\pi$, the polylogarithm terms disappear and the
identity becomes the compact evaluation
\begin{equation}\label{power:logspecial}
 \FP\int_0^1\gamma_0(x)
       \left\{\log\frac{\sin\pi x}{\pi}+\epsilon\pi ix\right\}\,\dd x
 =\gamma_1+\frac{\gamma^2}{2}+\frac{\pi^2}{24}
       -\frac{\log^2(2\pi)}2+\frac{\epsilon\pi i\log(2\pi)}2.
\end{equation}
The imaginary part is also an ordinary consequence of Raabe's
log-Gamma integral.  The constant real part supplies a useful
independent check of the branch and finite-part conventions.
The $q=0$ specialization is compatible with the prior sine-power
completion; the arbitrary nonreal $z$ identity is the new extension.

\subsection{Exact derivatives and normalized primitives in the resolvent parameter}

Denote the left side of \eqref{power:allmoments} by
$M_{m,p;L,k}(z)$.  Coefficientwise differentiation in $z$ is legitimate
by Theorem~\ref{power:completion}.  Since
$\partial_z\Log R_z=R_z$, we obtain the exact ladder
\begin{equation}\label{power:zladder}
 \frac{\dd}{\dd z}M_{m,p;L,k}(z)
    =L\,M_{m,p;L+1,k}(z)+k\,M_{m,p;L+1,k-1}(z),
\end{equation}
with the last term omitted for $k=0$.  This differentiates the actual
fixed-coordinate moments, rather than an uncompleted spectral limit.

For $L\ne1$ define
\begin{equation}\label{power:primitive}
 P_{m,p;L,k}(z)=
 \sum_{j=0}^k\frac{(-1)^jk!}{(k-j)!(L-1)^{j+1}}
                         M_{m,p;L-1,k-j}(z).
\end{equation}
For $L=1$ put $P_{m,p;1,k}=M_{m,p;0,k+1}/(k+1)$.
Equation \eqref{power:zladder} telescopes to $P'=M$.
Consequently
\begin{equation}\label{power:normalizedprimitive}
 \int_{z_0}^{z}M_{m,p;L,k}(\xi)\,\dd\xi
       =P_{m,p;L,k}(z)-P_{m,p;L,k}(z_0)
\end{equation}
for any path within the same open half-plane.  This fixes every additive
constant explicitly.  In particular \eqref{power:logexplicit} is an
explicit primitive of the ordinary $\gamma_0$ resolvent integral.

\paragraph{An exact boundary counterexample to naive substitution.}
The meromorphic transform is not, at every exceptional point, the
pointwise integral of the specialized integrand.  At $p=2$, $u=1$,
$\lambda=1$, the latter is zero because
$\partial_x^2\zeta(0,x)=0$.  Continuing \eqref{power:hurwitz} along
$\lambda=1$ first gives $-1$.  Indeed
$\mathcal L_1(u-2,q)=(q-1)\Li_{u-2}(q)/q$, with its removable
$q=0$ interpretation, gives $-1$ at $u=1$ by elementary algebra.
The sufficient convergence domain has $\Re(u+\lambda)>p$ and
excludes this boundary.  Locally the discrepancy is the limit of
$(u-1)(u-2)/(u-1)=u-2$.  This example identifies a concrete
normalization error in an uncompleted extension; it is not an allegation
that the inspected incoming resolvent theorem makes that error.

% END sections/complex_powers.tex


% BEGIN sections/moments.tex
\section{Exact ordinary moments of continuous resolvent powers}
\label{moment:sec}

The continuous-power transform has a useful specialization that remains
an ordinary convergent integral throughout its stated domain. It gives
every polynomial moment from a regularized Gauss function, and at the
two distinguished imaginary parameters the answer becomes a Gamma
quotient. Generalized harmonic numbers then organize all Bernoulli
moments at once.

For $z\in\mathbb C\setminus\mathbb R$, retain
\[
 \varepsilon=\operatorname{sgn}\operatorname{Im}z,\qquad
 C=-\frac1{z+\varepsilon\pi i},\qquad
 q=\frac{z-\varepsilon\pi i}{z+\varepsilon\pi i},\qquad
 \sigma=-2\varepsilon\pi i.
\]
Thus $|q|<1$. On $0<x<1$ put
\[
 R_z(x)=\frac1{\pi\cot(\pi x)-z},\qquad
 R_z(x)^\lambda=\exp\bigl(\lambda\operatorname{Log}R_z(x)\bigr),
\]
where the principal logarithm is continuous along the interval. Its
argument runs from $0$ to $\varepsilon\pi$. Powers of $C$ use its
principal logarithm as well. In particular,
\[
 R_z(x)^\lambda=C^\lambda
 \left(\frac{1-e^{-\sigma x}}{1-qe^{-\sigma x}}\right)^\lambda,
\]
where the power in parentheses is the radial boundary value of the
holomorphic power that equals one at the origin of the disk.

We use Olver's regularized Gauss function
\begin{equation}\label{moment:eq:regularized}
 \mathbf F(a,b;c;q)=
 \sum_{n=0}^\infty
 \frac{(a)_n(b)_n}{\Gamma(c+n)n!}q^n.
\end{equation}
When $c$ is not a nonpositive integer, this equals
${}_2F_1(a,b;c;q)/\Gamma(c)$. The series definition is essential at
apparent parameter resonances.

\begin{theorem}[Exponential moment generator]
\label{moment:thm:exponential}
For $\operatorname{Re}\lambda>-1$ and every $t\in\mathbb C$, set
$\xi=t/\sigma$. Then
\begin{equation}\label{moment:eq:exponential}
 \boxed{\displaystyle
 \int_0^1 e^{tx}R_z(x)^\lambda\,dx
 =C^\lambda e^{t/2}
   \frac{\Gamma(1+\lambda)}{\Gamma(1+\xi)}
   \mathbf F(\lambda,-\xi;1+\lambda-\xi;q).}
\end{equation}
Both sides are entire functions of $t$. They are holomorphic in
$\lambda$ throughout the half-plane $\operatorname{Re}\lambda>-1$.
\end{theorem}

\begin{proof}
The endpoint expansions are
\[
 R_z(x)=x(1+O(x)),\qquad
 R_z(1-x)=-x(1+O(x)).
\]
The prescribed branches therefore bound the integrand, uniformly on
compact parameter sets, by an integrable constant times
$[x(1-x)]^{\alpha}$, with $\alpha>-1$. Parameter derivatives add
only powers of endpoint logarithms. This proves convergence and the
stated parameter holomorphy for the integral.

Write
\[
 f(w)=\left(\frac{1-w}{1-qw}\right)^\lambda
      =\sum_{n\ge0}d_nw^n,
 \qquad d_0=1.
\]
First suppose $\operatorname{Re}\lambda>0$ and
$\operatorname{Re}\xi<0$. The coefficients are absolutely summable:
the coefficients of $(1-w)^\lambda$ are
$O(n^{-1-\operatorname{Re}\lambda})$, with the terminating cases
understood separately, while $(1-qw)^{-\lambda}$ has geometrically
decaying coefficients. Consequently termwise integration gives
\[
 C^{-\lambda}\int_0^1e^{tx}R_z(x)^\lambda\,dx
 =(e^t-1)\sum_{n\ge0}\frac{d_n}{t-\sigma n}
 =-\frac{e^t-1}{\sigma}\int_0^1
        u^{-\xi-1}(1-u)^\lambda(1-qu)^{-\lambda}\,du.
\]
Euler's integral for the regularized Gauss function
\cite[15.6.1]{DLMF} evaluates the last expression as
\[
 \frac{e^t-1}{t}
 \Gamma(1-\xi)\Gamma(1+\lambda)
 \mathbf F(\lambda,-\xi;1+\lambda-\xi;q).
\]
The Gamma reflection formula and $\sigma=-2\varepsilon\pi i$ give
\[
 \frac{e^t-1}{t}\Gamma(1-\xi)
       =\frac{e^{t/2}}{\Gamma(1+\xi)}.
\]
This proves \eqref{moment:eq:exponential} on the initial open domain.

For $|q|<1$, the series \eqref{moment:eq:regularized} converges
normally on compact subsets of its parameters and is entire in its
three parameters. The right side of \eqref{moment:eq:exponential}
is therefore entire in $t$ and holomorphic for
$\operatorname{Re}\lambda>-1$. The identity theorem first extends
in $t$ and then in $\lambda$, proving the full statement. In
particular, no exclusion of $1+\lambda-\xi=0,-1,-2,\ldots$ is
necessary in the regularized notation.
\end{proof}

The ingredients in this proof are classical Beta and Gauss identities.
The contribution here is their branch-consistent specialization to the
continuous-power resolvent and its exact moment consequences.

\begin{corollary}[Means and logarithmic moments]
\label{moment:cor:mean}
For $\operatorname{Re}\lambda>-1$ and $k\ge0$,
\begin{equation}\label{moment:eq:meanlogs}
 \int_0^1 R_z(x)^\lambda
       \bigl(\operatorname{Log}R_z(x)\bigr)^k\,dx
       =C^\lambda(\operatorname{Log}C)^k.
\end{equation}
In particular, $\int_0^1R_z(x)^\lambda\,dx=C^\lambda$ and
$\int_0^1(\operatorname{Log}R_z(x))^k\,dx=(\operatorname{Log}C)^k$.
\end{corollary}
\begin{proof}
Put $t=0$ in \eqref{moment:eq:exponential}, using
$\mathbf F(\lambda,0;1+\lambda;q)=1/\Gamma(1+\lambda)$.
Differentiate $k$ times in $\lambda$. The endpoint majorants in the
proof of the theorem justify each differentiation under the integral.
\end{proof}

\subsection{The Bernoulli generator and its Gamma specialization}

Let $B_m(x)$ be the Bernoulli polynomials with
$B_1(x)=x-1/2$. Multiplication of
\eqref{moment:eq:exponential} by $t/(e^t-1)$ gives the exact generator
\begin{align}\label{moment:eq:bernoulli-general}
 \sum_{m\ge0}\frac{t^m}{m!}
   \int_0^1 B_m(x)R_z(x)^\lambda\,dx
 ={}&C^\lambda\Gamma(1-t/\sigma)\Gamma(1+\lambda)\notag\\
 &\quad\times
 \mathbf F(\lambda,-t/\sigma;1+\lambda-t/\sigma;q),
 \qquad |t|<2\pi.
\end{align}
The Bernoulli series converges uniformly in $x\in[0,1]$ on every
smaller $t$-disk. The same integrable endpoint majorant justifies the
interchange with the integral. Thus coefficient extraction is an
ordinary-integral identity, including when $\lambda=-1/2$.

\begin{corollary}[A digamma--Lerch first moment]
\label{moment:cor:first}
For $\operatorname{Re}\lambda>-1$,
\begin{align}\label{moment:eq:first}
 \int_0^1\left(x-\frac12\right)R_z(x)^\lambda\,dx
 =\frac{C^\lambda}{\sigma}\bigl[&\psi(1+\lambda)+\gamma
    +q\Phi(q,1,1+\lambda)\notag\\
    &+\operatorname{Log}(1-q)\bigr],
\end{align}
where $\Phi(q,1,a)=\sum_{n\ge0}q^n/(n+a)$.
\end{corollary}
\begin{proof}
With $f(w)$ as above, its nonconstant Fourier coefficients give
\[
 C^{-\lambda}\int_0^1\left(x-\frac12\right)R_z(x)^\lambda dx
       =-\frac1\sigma\int_0^1\frac{f(v)-1}{v}\,dv.
\]
This follows first for $\operatorname{Re}\lambda>0$ by absolute
summation and then for $\operatorname{Re}\lambda>-1$ by holomorphy.
For $0<q<1$, substitute $u=(1-v)/(1-qv)$ in the last integral:
\[
 \int_0^1\frac{f(v)-1}{v}\,dv
 =\int_0^1(u^\lambda-1)
       \left(\frac1{1-u}-\frac q{1-qu}\right)du
 =-\psi(1+\lambda)-\gamma-q\Phi(q,1,1+\lambda)
      -\log(1-q).
\]
The first equality extends to $|q|<1$ by holomorphy. Substitution
proves \eqref{moment:eq:first} with the principal logarithm.
\end{proof}

Define complex generalized harmonic numbers on
$\operatorname{Re}\lambda>-1$ by
\begin{equation}\label{moment:eq:harmonics}
 H_\lambda^{(1)}=\psi(1+\lambda)+\gamma,
 \qquad
 H_\lambda^{(j)}=\zeta(j)-\zeta(j,1+\lambda)\quad(j\ge2).
\end{equation}
For nonnegative integers $\lambda=N$ these are the usual finite
harmonic numbers. Define $h_m(\lambda)$ by
\begin{equation}\label{moment:eq:hgenerator}
 \sum_{m\ge0}h_m(\lambda)u^m
 =\frac{\Gamma(1-u)\Gamma(1+\lambda)}
        {\Gamma(1+\lambda-u)}
 =\exp\left(\sum_{j\ge1}\frac{H_\lambda^{(j)}}j u^j\right)
 \qquad (|u|<\min\{1,|1+\lambda|\}).
\end{equation}
The Taylor series of the Gamma quotient itself is holomorphic for
$|u|<1$. The smaller disk in the exponential representation accounts
for possible zeros of that quotient when $\lambda$ has negative real
part; its logarithmic series must not be continued through such a zero.
Equivalently,
\begin{equation}\label{moment:eq:hrecurrence}
 h_0=1,\qquad
 h_m=\frac1m\sum_{j=1}^m H_\lambda^{(j)}h_{m-j}\quad(m\ge1).
\end{equation}
These are complete homogeneous polynomials in the generalized
harmonic coordinates; for example,
\[
 h_1=H_\lambda^{(1)},\quad
 h_2=\frac{(H_\lambda^{(1)})^2+H_\lambda^{(2)}}2,\quad
 h_3=\frac{(H_\lambda^{(1)})^3+
        3H_\lambda^{(1)}H_\lambda^{(2)}+2H_\lambda^{(3)}}6.
\]

\begin{theorem}[Every Bernoulli moment at the imaginary base points]
\label{moment:thm:bernoulli}
For $z=\varepsilon\pi i$, $\operatorname{Re}\lambda>-1$, and
every integer $m\ge0$,
\begin{equation}\label{moment:eq:allmoments}
 \boxed{\displaystyle
 \int_0^1 B_m(x)
  \left(\frac{\sin(\pi x)}\pi e^{\varepsilon\pi i x}\right)^\lambda dx
 =m!\left(\frac{\varepsilon i}{2\pi}\right)^\lambda
       (-2\varepsilon\pi i)^{-m}h_m(\lambda).}
\end{equation}
The power on the left means
$\exp\{\lambda[\log(\sin(\pi x)/\pi)+\varepsilon\pi i x]\}$.
\end{theorem}
\begin{proof}
At these points $q=0$ and $C=\varepsilon i/(2\pi)$, while
$R_{\varepsilon\pi i}(x)=\sin(\pi x)e^{\varepsilon\pi i x}/\pi$.
Equation \eqref{moment:eq:bernoulli-general} therefore reduces to
$C^\lambda$ times the Gamma quotient
\eqref{moment:eq:hgenerator} at $u=t/\sigma$. Comparing coefficients
proves the formula. The logarithmic Taylor expansion of that quotient
gives \eqref{moment:eq:hgenerator} from
\eqref{moment:eq:harmonics}; logarithmic differentiation then gives
\eqref{moment:eq:hrecurrence}.
\end{proof}

At the negative half power,
\[
 H_{-1/2}^{(1)}=-2\log2,\qquad
 H_{-1/2}^{(j)}=-(2^j-2)\zeta(j)\quad(j\ge2).
\]
The first two nonconstant moments are consequently
\begin{align}\label{moment:eq:halfexamples}
 \int_0^1\left(x-\frac12\right)
  \sqrt{\frac\pi{\sin(\pi x)}}e^{-\varepsilon\pi i x/2}\,dx
 &=-\frac{1+\varepsilon i}{\sqrt\pi}\log2,\notag\\
 \int_0^1\left(x^2-x+\frac16\right)
  \sqrt{\frac\pi{\sin(\pi x)}}e^{-\varepsilon\pi i x/2}\,dx
 &=\sqrt\pi(1-\varepsilon i)
     \left(\frac1{12}-\frac{\log^22}{\pi^2}\right).
\end{align}
Both endpoints have only inverse-square-root singularities, so these
are absolutely convergent ordinary integrals. For example the second
identity also yields the entirely real evaluation
\begin{equation}\label{moment:eq:realhalf}
 \int_0^1\left(x^2-x+\frac16\right)
 \frac{\cos(\pi x/2)}{\sqrt{\sin(\pi x)}}\,dx
 =\frac1{12}-\frac{\log^22}{\pi^2}.
\end{equation}
At the positive half power,
$H_{1/2}^{(1)}=2-2\log2$ and
$H_{1/2}^{(j)}=2^j-(2^j-2)\zeta(j)$ for $j\ge2$.
Thus every moment at either half power belongs to an explicitly
specified algebra of zeta values, $\log2$, $\pi$, and the displayed
square-root prefactor. Further $\lambda$-derivatives of
\eqref{moment:eq:allmoments} give all corresponding moments with
arbitrary powers of $\operatorname{Log}R_{\varepsilon\pi i}$.

\paragraph{Verification.}
The independent script \texttt{check\_moments.py} compares the
Gamma--Gauss generator against direct real-interval quadrature in both
half-planes, including nonintegral powers with negative real part and
powers with real part greater than one. It also compares four Bernoulli
moments with \eqref{moment:eq:allmoments}, reconstructs the first nine
Fourier coefficients by two exact symbolic formulas, and rejects a
changed coefficient. Its separate spectral and finite-part checks audit
the conventions used in the surrounding continuous-power transform.
These numerical comparisons are diagnostics, not interval certificates
or steps in the proofs above.

% END sections/moments.tex


% BEGIN sections/audit_and_integration.tex
\section{Source audit and proposed integration}
\label{audit:section}

\subsection{What the audit established}

No new false theorem was found in the specific source arguments used
for the present results. In particular, the incoming cubic bridge,
the centered parity criterion, and the previously proved
trigamma-square moments survive the checks relevant to their use
here. Their stated limits are research questions, not flaws in
their proofs. This is a targeted mathematical audit of the relevant
material, not a certification of every page in the repository.

The appropriate source changes are therefore progress updates and
additional normalization examples. Table~\ref{audit:status-table}
gives the precise mapping. The supplied integration notes identify
the archive-internal source paths and theorem labels; no source
repository was modified in preparing this package.

\begin{table}[htbp]
\centering\small
\begin{tabular}{@{}>{\raggedright\arraybackslash}p{.28\linewidth}>{\raggedright\arraybackslash}p{.36\linewidth}>{\raggedright\arraybackslash}p{.28\linewidth}@{}}
\toprule
Source question & Result established here & Remaining scope\\
\midrule
Cubic harmonic report, Q6 & Quartic power-to-harmonic reduction,
Theorem~\ref{q4:main} & Powers $k\ge5$ and arithmetic reduction of
the retained depth-three coefficient\\[5pt]
Harmonic-parity report, R1 & All pure two-even-tail moments and their
generator, Theorems~\ref{tail:all-moments}, \ref{tail:generator}
& Additional harmonic factors and products of four or more tails\\[5pt]
Harmonic-parity report, R4 & Complex powers, joint resonant completion,
all logarithmic moments, Theorem~\ref{power:completion}
& Products of distinct resolvents or nonlinear continuous-order factors\\[5pt]
Higher-reflection report, centered products & First mixed pair and all
two-tail orders & Higher-depth products and spectral derivatives\\[5pt]
Canonical Gaussian $S_6$, revised $S_8$ & No proof or refutation obtained
& Exact functional bridge to the frozen candidates\\
\bottomrule
\end{tabular}
\caption{Scoped updates relative to the pinned source revision.}
\label{audit:status-table}
\end{table}

\subsection{A concrete correction to a tempting extension}

There is a mathematically incorrect extension that the new results
rule out: replacing a specialized coordinate finite part by the
meromorphic transform at every point where the specialized integrand
happens to be integrable. The counterexample at the end of
Section~\ref{power:sec} is exact. With $p=2$, $u=1$, and
$\lambda=1$, the ordinary integral of
$\partial_x^2\zeta(0,x)R_z(x)$ equals zero. The fixed-$\lambda$
meromorphic continuation of the spectral formula equals $-1$.
The local monomial has a coefficient that vanishes at the same
point as its integration denominator, leaving the nonzero limit
$u-2$.

The correction is to retain the entire moving amplitude before
coefficient extraction, as in \eqref{power:H}. Removing just the
negative Laurent coefficients of a one-variable expansion is
insufficient. At negative integer powers the two smooth endpoint
residues cancel, yet their quotient has the nonzero limit
\eqref{power:smoothcontact}. A cancellation of residues therefore
does not imply that there is no finite endpoint correction.

This is a warning and an exact regression example for future
extensions. The inspected incoming resolvent theorem states its
own convergence and completion conditions; this article does not
attribute the incorrect extension to that theorem. The package
includes a deliberately corrupted coefficient check so that
omitting the endpoint phase is visibly rejected by exact algebra.

\subsection{Four conventions that should accompany integration}

\begin{enumerate}
\item \textbf{Keep strict harmonic coordinates named and normalized.}
The quartic coefficient $\theta$ is the coefficient of $u$ in
$Z_3(1+u,1,1)$, not the constant coefficient and not an order
derivative of a regular function before its poles have been
specified. The entire principal part and constant coefficient
are given in \eqref{q4:theta}. The retained coefficient is
genuine remaining arithmetic data in the present reduction;
its integral representation is a normalization, not an ordinary-
constant evaluation.

\item \textbf{Retain subtractions inside convergent summands.}
The sums in \eqref{tail:ordinary-formula} and
\eqref{tail:rational-polygamma} are ordinary convergent sums of
the displayed brackets. Their polynomial pieces must not be
separated and summed as ordinary series. The equality with the
continued Dirichlet series is proved using the half-shift and
Hurwitz trivial zeros.

\item \textbf{Use the regularized Gauss function at parameter resonances.}
The definition \eqref{moment:eq:regularized} remains valid when
its third parameter is a nonpositive integer. A quotient written
informally as ${}_2F_1/\Gamma(c)$ must be read through that
definition. Likewise the logarithmic Gamma-quotient series has
the smaller disk stated in \eqref{moment:eq:hgenerator}; zeros
of the quotient obstruct its logarithm before they obstruct
the quotient itself.

\item \textbf{Distinguish analytic proofs, finite algebra, and numerical evidence.}
The infinite identities follow from the proofs. Symbolic checks
verify selected finite polynomial consequences. The numerical
files use ordinary multiprecision arithmetic with specified
analytic tail estimates. They are not interval-enclosed evaluations
and do not certify arithmetic independence of any constants.
\end{enumerate}

\subsection{The Gaussian conjectures and the exact obstruction}

The frozen $S_6$ target is in
\texttt{chapters/04-cyclotomic-quotients.tex}; the revised $S_8$
candidate is in \texttt{chapters/04-S8-candidate.tex} of the
manuscript at the pinned revision~\cite{ProveIt}. Both remain
conjectures. Their existing short candidate formulas have not
been replaced by a numerical fit in this work.

The exact-resonant report already proves a specific obstruction:
the $S_6$ target is separated from the entire convergent Cayley
shuffle ideal together with a specified family of $5{,}131$
additional relations~\cite{IncomingExact}. This is stronger than
an unsuccessful small linear solve. It is also a statement about
that named formal relation system, not a proof that the numerical
period conjecture is false or that no other functional relation
can prove it.

The larger search in the independent-orders report generated
$51{,}244$ double-shuffle rows in $24{,}514$ imaginary word
coordinates. Its modular elimination stopped after $8{,}500$
pivots, before exhausting the system~\cite{IncomingIndependent}.
The incomplete search cannot be cited as an exhausted obstruction,
and it supplies no rational identity certificate. No new claim
about its terminal rank is made here.

The new complex-power transform provides analytic identities that
can be tested against a Gaussian target, but this article has not
proved a specialization producing either frozen residual. The
unresolved step is an explicit functional bridge, followed by
exact coefficient extraction. Giving another representation of
an unrelated zeta constant does not accomplish that step.


% END sections/audit_and_integration.tex


% BEGIN sections/research_questions.tex
\section{Further research questions and an order of attack}
\label{research:section}

The following questions ask for exact identities, finite reductions,
or normalized analytic representations. They are not asserted
conjectures unless an implication is explicitly proved above.
The requested deliverable for each question is stated to separate
mathematical progress from additional numerical evidence.

\subsection{Higher digamma powers and harmonic Laurent coefficients}

\paragraph{1. The fifth power with all finite corrections.}
Repeat the quartic contour and finite summation argument for
$\FP\int_0^1\psi(x)^5\,\dd x$. The principal parts are
algorithmically finite harmonic polynomials, but their integration
requires a new exact summation-by-parts reduction. Determine
which linear Laurent coefficients of strict depth-four and
lower-depth zeta functions enter, and give every ordinary zeta
derivative and Stieltjes correction. A proof should include the
constant term of each finite asymptotic sum, not only matching
principal parts or a high-precision decimal.

\paragraph{2. Canonical polynomials that remove lower-depth data.}
The polynomial $T^4+4\gamma T^3$ removes $\eta$ from the quartic
bridge. Is there a systematically constructed monic polynomial
$P_k(T)$, with coefficients in a stated ordinary-constant algebra,
for which $\FP\int_0^1P_k(\psi(x))\,\dd x$ retains only one
new highest-depth harmonic coefficient? The Gamma and inverse-Gamma
coefficient polynomials are natural candidates for organizing the
recursion. A useful theorem would give the recurrence and prove
the cancellation for arbitrary $k$; the single quartic cancellation
does not establish such a theorem.

\paragraph{3. The arithmetic content of the depth-three coordinate.}
The normalized elementary kernel \eqref{q4:theta-kernel} determines
$\theta$ exactly. Seek an independent identity placing it in a
smaller algebra of zeta derivatives, ordinary Stieltjes constants,
or Gamma/Barnes integrals. The target should state which coordinates
are eliminated. Merely replacing the strict Laurent definition
by another name for the same integral does not reduce the
arithmetic content. Conversely, a failed numerical relation search
does not prove irreducibility.

\paragraph{4. A common-Hurwitz-shift quartic bridge.}
For $a>0$, define a consistent common shift in every denominator
of the strict depth-three series and derive the counterpart of
\eqref{q4:Didentity}. Determine the associated identity for
products of $\psi(a+x)$ and its derivatives, including the two
ordinary endpoint values on $[a,a+1]$. Then analyze the singular
limit $a\downarrow0$ with a specified coordinate finite part.
Replacing $\gamma_m$ by $\gamma_m(a)$ in the unshifted formula
without deriving these endpoint terms is not valid evidence.

\subsection{Harmonic-tail arithmetic beyond two factors}

\paragraph{5. The remaining harmonic factors in R1.}
For the products
\[
A_p(n)A_q(n)H_n^a
\prod_{j\ge1}\bigl(H_n^{(2j+1)}\bigr)^{b_j},
\]
with finitely many nonzero $b_j$, determine finite formulas for
the regular centered values and a moment-independent coordinate
space for each fixed polynomial. The source parity theorem predicts
local regularity, but it does not evaluate the global finite
constant. Finite summation by parts increases the depth of the
harmonic sums. One must specify the necessary multiple-zeta
coordinates and prove any claimed depth reduction.

\paragraph{6. Four even tails and the first genuine depth increase.}
Carry out the analogous calculation for
$A_pA_qA_rA_s$, first for $A_2^4$ and $A_2^3A_4$.
An even number of even tails preserves the centered even-power
asymptotics. The two-tail theorem suggests a finite Bernoulli
reduction at every moment order, but higher-depth convergent sums
can survive. Determine a fixed spanning set and an analytic
generating function. A parity argument alone does not prove that
all such higher-depth values reduce to ordinary zeta products.

\paragraph{7. First spectral derivatives at the regular lattice.}
Find exact reductions of $D'_{p,q}(-2r)$ or combinations of these
derivatives. Differentiating the subtracted continuation produces
logarithmically weighted global remainder sums as well as
derivatives of the explicitly removed Hurwitz terms. Seek
identities of these kernels that cancel the global terms, or
identify a minimal additional harmonic-Stieltjes coordinate
family. The finite value formula cannot simply be differentiated
with respect to its discrete integer label $r$.

\paragraph{8. Finite generation by initial moments.}
For fixed even $p,q$, Corollary~\ref{tail:fixed-space} gives at
most $\max(p,q)$ specified arithmetic coordinates. Determine
whether a uniform finite prefix of the moment sequence generates
the same rational span, and exhibit a rational inverse matrix when
it does. The $(2,4)$ case is proved by its first four moments.
An optimal bound concerns the rank of a specified coefficient
matrix. Equality with a dimension of actual constants would require
an additional arithmetic independence result.

\paragraph{9. Rational shifts and colored polylogarithms.}
Replace the inner shift $n+1$ and the center $n+1/2$ by compatible
rational shifts. Classify when the removed polynomial sums vanish,
and determine the surviving finite corrections when they do not.
The natural targets are Hurwitz zeta values at rational arguments
and colored Euler sums, together with polylogarithms at roots of
unity. A paired formulation at shifts $a$ and $1-a$ may retain
more symmetry than an isolated noncentral shift.

\subsection{Resolvent products, spectral derivatives, and primitives}

\paragraph{10. Several distinct nonreal resolvents.}
For
\[
\prod_{j=1}^d R_{z_j}(x)^{\lambda_j},
\]
derive an explicit joint Fourier or beta transform and complete
the hyperplanes where the total endpoint exponent is integral.
When all $z_j$ lie in the same half-plane the disk factors have
a common Fourier orientation; an Appell or Lauricella function
is a plausible ordinary-moment generator. Parameters in opposite
half-planes introduce two-sided convolutions and additional
arithmetic structure. Each claim needs a branch prescription and
the full endpoint amplitude before any specialization.

\paragraph{11. Products of continuous polygamma orders.}
The transform in this article is linear in the Hurwitz factor.
For products of two independently varying Hurwitz orders, derive
the convolution and identify exactly when ordinary polylogarithm
order derivatives suffice and when ordered or Tornheim coordinates
remain. Mixed order derivatives should be extracted from a joint
completed germ. Multiplying two separately completed scalar
identities does not account for the product's local singularities.

\paragraph{12. Arithmetic at rational complex-power exponents.}
At $q=0$, every Bernoulli moment is a finite polynomial in the
generalized harmonic values $H_\lambda^{(j)}$ and the fixed
prefactor. At rational $\lambda$ these values have classical
Hurwitz-zeta descriptions. Determine comparable finite reductions
for useful nonzero algebraic $q$ and rational $\lambda$, especially
for derivatives in $\lambda$ and the polylogarithm order.
The beta--polylogarithm formula supplies a convergent starting
integral; it does not by itself prove a finite constant evaluation.

\paragraph{13. Normalized primitives beyond the balanced combinations.}
Combine the explicit $z$-ladder with argument integration and
Hurwitz order differentiation to seek individual Barnes-type
representations. State the additive normalization at a base point
and derive its transport under $z$ and argument shifts. The
existing difference primitives already fix their constants; a
new result should identify a smaller named function class or an
additional functional equation, rather than introduce an arbitrary
integration constant.

\subsection{The Gaussian bridge and exact certification}

\paragraph{14. A functional identity for the frozen Gaussian residual.}
Express the actual $S_6$ or revised $S_8$ residual as a coefficient
of a proved joint generating identity, with all changes of
variables and branches explicit. The Cayley-ideal obstruction
specifies information that its named relation family cannot
supply. A successful new relation should have a computable
nonzero action on the separating functional, or an analytic
specialization demonstrably outside that family. The article's
resolvent transform is a source of candidate relations, not an
already established bridge.

\paragraph{15. Complete and certify any enlarged relation search.}
If the larger modular calculation is resumed, first preserve the
frozen target, relation enumeration, and checkpoint hashes. Finish
the declared relation system and recover an exact rational
certificate or exact separating functional. A modular rank
observation alone is not a rational proof of the target identity,
and an interrupted elimination does not show nonmembership in
the generated span. The input system and the terminal conclusion
must be independently replayable.

\paragraph{16. Formalize reusable analytic and algebraic components.}
The most reusable formal units are the two-point
Mittag--Leffler subtraction, the finite Faulhaber constant-term
calculation, the beta coefficient identity
\eqref{power:beta-coeff}, and the moving-monomial completion lemma
inside Theorem~\ref{power:completion}. Formalizing these would
make subsequent all-index identities considerably easier to
verify. An interval evaluation of $\theta$ or $Q_4$ can be added
separately by enclosing the same analytic tail formulas; it
would strengthen the numerical artifact, not replace the exact
proof.

\subsection{Practical sequence}

The most direct next calculations are the fifth-power reduction,
the initial-moment coefficient matrices, and the first additional
harmonic factors in the two-tail family. Each has a finite exact
starting point supplied here. Four-tail products and joint
resolvent powers are natural next structural extensions. The
Gaussian reduction should proceed only with a concrete candidate
functional bridge or a completed exact relation computation;
more precision for the same unproved scalar fit does not remove
the known missing step.


% END sections/research_questions.tex

\appendix

% BEGIN sections/mixed_even_tail_verification.tex
\section{Verification of the mixed-tail identities}
\label{tail:verification-section}

The script \texttt{verification/verify\_mixed\_tails.py} records
finite exact checks and a numerical route independent of the moment
formula. Its numerical tail approximation has an explicit analytic
remainder estimate. Ordinary multiprecision arithmetic is used, so
the recorded floating-point computations are not asserted to be
interval certificates. The identities are established by the
proofs above.

\subsection{A proved truncation envelope}

For real \(x>0\), put
\[
 T_p(x)=\zeta(p,x+1/2),\qquad
 Q_{p,K}(x)=\sum_{j=0}^{K}c_{p,j}x^{1-p-2j},\qquad
 R_{p,K}=\frac{2(p+2K)!}{(p-1)!(2\pi)^{2K+2}}.
\]
Then
\begin{equation}\label{tail:individual-bound}
 0<T_p(x)\le\frac{x^{1-p}}{p-1},\qquad
 |T_p(x)-Q_{p,K}(x)|\le R_{p,K}x^{-p-2K-1}.
\end{equation}
Here and below the bound is needed only to substantiate numerical
verification, not as the research objective.

To prove it, let \(g(u)=u/(2\sinh(u/2))\). The standard
hyperbolic-cosecant partial fractions \cite[Eq.~4.36.5]{DLMF}
give
\[
 g(u)=1+2\sum_{m=1}^{\infty}
             (-1)^m\frac{u^2}{u^2+(2\pi m)^2}.
\]
Expanding each rational summand through degree \(2K\), and using
the alternating-series estimate on the remainder, gives for real
\(u>0\)
\[
 \left|g(u)-\sum_{j=0}^{K}
        \frac{B_{2j}(1/2)}{(2j)!}u^{2j}\right|
 \le\frac{2u^{2K+2}}{(2\pi)^{2K+2}}.
\]
The identity
\[
 T_p(x)=\frac1{(p-1)!}\int_0^{\infty}
             u^{p-2}e^{-xu}g(u)\,\dd u
\]
now proves both parts of \eqref{tail:individual-bound}; the
positivity bound uses \(0<g(u)\le1\).

Fix \(N\ge1\), put \(a=N+1/2\), and choose \(K\ge K_r\).
Define
\[
 M_{p,K}(a)=\sum_{j=0}^{K}|c_{p,j}|a^{-2j},\qquad
 k_{p,q,\ell}^{[K]}
   =\sum_{\substack{0\le j\le K\\0\le\ell-j\le K}}
                         c_{p,j}c_{q,\ell-j}.
\]
The finite computational expression is
\begin{align}\label{tail:finite-evaluation}
 V_{p,q,r}^{N,K}={}&
 \sum_{n=0}^{N-1}\left\{x_n^{2r}A_p(n)A_q(n)
     -\sum_{\ell=0}^{K_r-1}
            k_{p,q,\ell}x_n^{2r+2-W-2\ell}\right\}\notag\\
 &+\sum_{\ell=K_r}^{2K}
       k_{p,q,\ell}^{[K]}\zeta(W-2+2\ell-2r,a).
\end{align}
All the Hurwitz-zeta orders in the last sum are at least two.
Provided \(W+2K-2r>1\), its error satisfies
\begin{align}\label{tail:evaluation-envelope}
 |D_{p,q}(-2r)-V_{p,q,r}^{N,K}|
 \le{}&\left\{\frac{R_{p,K}}{q-1}
                   +M_{p,K}(a)R_{q,K}\right\}
                     \zeta(W+2K-2r,a).
\end{align}
One may take the smaller bound obtained by interchanging \(p,q\).
Indeed, \eqref{tail:individual-bound} gives
\[
 |T_pT_q-Q_{p,K}Q_{q,K}|
 \le\left\{\frac{R_{p,K}}{q-1}
                   +M_{p,K}(a)R_{q,K}\right\}x^{-W-2K}
 \quad(x\ge a).
\]
Multiplication by \(x^{2r}\), summation over the omitted tail,
and expansion of the finite product prove
\eqref{tail:evaluation-envelope}.

\subsection{Recorded checks}

The script uses \(110\) decimal working digits, \(N=64\), and
\(K=30\). It checks \(r=0,1,2,3,5,6\) for each of
\[
 (p,q)=(2,2),(2,4),(4,4),(2,6),(4,6),(6,6).
\]
These include the first regularized values and several higher
regularized values, not only ordinary unsubtracted sums.
Every one of the \(36\) observed discrepancies is below its
analytic truncation envelope. Selected results are
\begin{center}
\begin{tabular}{cccll}
\hline
\(p\)&\(q\)&\(r\)&Absolute discrepancy&Truncation envelope\\
\hline
2&4&2&\(1.4110\,10^{-73}\)&\(1.4471\,10^{-73}\)\\
2&4&3&\(5.9899\,10^{-70}\)&\(6.1431\,10^{-70}\)\\
2&4&6&\(4.6138\,10^{-59}\)&\(4.7314\,10^{-59}\)\\
4&4&6&\(7.2161\,10^{-63}\)&\(7.4014\,10^{-63}\)\\
4&6&6&\(1.8213\,10^{-64}\)&\(1.8709\,10^{-64}\)\\
6&6&6&\(5.1274\,10^{-68}\)&\(5.2690\,10^{-68}\)\\
\hline
\end{tabular}
\end{center}

The exact checks independently expand the finite uncentered tail
products and Faulhaber polynomials and extract their constant
coefficients for six pairs and eleven odd block orders. They also
verify the short \((2,4)\) block formula through \(j=12\) and
recover the previously proved trigamma-square formula through
\(r=12\). The two eliminations in
\eqref{tail:24-eliminations} are checked by exact rational algebra.

Finally, direct quadrature of the generator is compared with its
first \(41\) Taylor terms for the three cases
\[
 (p,q,t)=(2,4,0.7),\quad (2,4,1.3),\quad (4,6,0.8).
\]
Their respective differences are approximately
\(1.1\,10^{-78}\), \(1.2\,10^{-56}\), and
\(2.6\,10^{-76}\), consistent with the omitted convergent Taylor
tail. For both \((2,4)\) checks, direct integral evaluation and
the explicit \(\operatorname{Li}_2\)--\(\operatorname{Li}_6\)
formula agree to better than \(10^{-109}\). Complete numerical
values, software versions, formulas, and truncation envelopes are
recorded in \texttt{verification/mixed\_tail\_results.json}.


% END sections/mixed_even_tail_verification.tex


% BEGIN sections/reproducibility.tex
\section{Reproducibility and verification record}
\label{repro:section}

The package contains the complete modular source, a self-contained
single-file TeX version, the compiled PDF, four verification
programs, their recorded JSON outputs, source provenance, a claim
inventory, and integration notes. It does not require a checkout
of ProveIt to build or run its mathematical checks. The source
snapshot is recorded for comparison and attribution.

\subsection{Build and replay}

From the extracted package directory, run
\begin{verbatim}
python3 build.py
python3 run_verification.py
\end{verbatim}
The first command generates the standalone TeX file and builds the
PDF using pdfLaTeX. The second runs the four mathematical suites
and saves a replay summary and per-suite logs under
\texttt{verification/replay/}. Each suite also updates its own
recorded JSON file. To preserve the shipped records verbatim,
run the replay on a copy of the package.

The Python dependencies are SymPy and mpmath, with the recorded
versions pinned in \texttt{requirements.txt}. The TeX build uses
standard AMS, Latin Modern, geometry, microtype, hyperref, and
related standard packages listed explicitly in the preamble.
No external bibliographic processor, private dataset, network
service, or shell-escape operation is needed.

\subsection{The scope of the checks}

The recorded exact checks total $183$, together with $21$
deliberate corruption controls. They are finite symbolic tests
of separate coefficient constructions and finite algebraic
consequences. They are not a count of independent theorems, and
they do not replace the all-index analytic arguments.

\begin{center}\small
\begin{tabular}{@{}>{\raggedright\arraybackslash}p{.27\linewidth}>{\raggedleft\arraybackslash}p{.13\linewidth}>{\raggedright\arraybackslash}p{.53\linewidth}@{}}
\toprule
Suite & Exact checks & Independent numerical route\\
\midrule
Quartic bridge & 44 & Harmonic/Binet summation versus local Taylor
integration and ordinary quadrature, at 100 working digits\\[4pt]
Mixed tails & 94 & 36 finite-head/Hurwitz-tail evaluations,
plus generator comparisons, at 110 working digits\\[4pt]
Continuous-power moments & 9 & 19 quadrature, Mellin, Gamma,
Lerch, and subtracted logarithmic comparisons, at 55 working digits\\[4pt]
Endpoint contacts and primitives & 36 & Exact arithmetic only;
20 omitted-phase controls are rejected\\
\bottomrule
\end{tabular}
\end{center}

The remaining corruption control changes a Fourier coefficient's
prefactor. It is rejected by exact polynomial comparison.
The continuous-power numerical suite samples both half-planes,
complex exponents, an integrable negative half power, and nonzero
complex $q$. Its largest absolute discrepancy is about
$2.824\cdot10^{-31}$, arising in inverse-square-root endpoint
quadrature. This is consistent with the deliberately more modest
diagnostic tolerance; it is not a claim of 55 correct digits in
every quadrature.

For the quartic calculation, the finest arithmetic evaluation and
the independent local Taylor/quadrature evaluation differ by less
than $1.40\cdot10^{-83}$. The Binet truncation bound is below
$1.88\cdot10^{-83}$. For the mixed tails, every observed moment
discrepancy is below its displayed analytic truncation envelope.
These statements separate the analytically controlled omitted tail
from floating-point rounding, which is not enclosed by this
implementation. No interval certificate is claimed.


% BEGIN sections/contact_verification.tex
\subsection{Exact contact and antiderivative checks}
\label{contact:verification}

The fixed-coordinate correction at a negative integral power has an
elementary residue interpretation.  For $M\ge1$,
\begin{equation}\label{contact:residue}
 [x^{M-1}]\{x(K(x)-z)\}^M
 =\operatorname*{Res}_{x=0}(K(x)-z)^M\,\dd x
 =\frac{(-z+i\pi)^M-(-z-i\pi)^M}{2i\pi}.
\end{equation}
Indeed, the local coordinate $v=1/K(x)=\tan(\pi x)/\pi$ satisfies
$\dd x=\dd v/(1+\pi^2v^2)$ and has derivative $1$ at zero.  Thus
the residue is the coefficient of $v^{-1}$ in
$(v^{-1}-z)^M/(1+\pi^2v^2)$, namely
\[
 \sum_{j=0}^{\lfloor(M-1)/2\rfloor}
   \binom{M}{2j+1}(-z)^{M-2j-1}(-\pi^2)^j.
\]
The binomial theorem gives the final expression in
\eqref{contact:residue}.  Consequently the difference between the
coordinate finite-part mean and its continued spectral value is
\begin{equation}\label{contact:identity}
 \frac{(-z+i\pi)^M+(-z-i\pi)^M}{2}
       -(-z-\epsilon i\pi)^M
 =\epsilon\pi i\,[x^{M-1}]\{x(K(x)-z)\}^M.
\end{equation}
This proves the contact identity for every $M$, independently of the
general completion theorem.

The accompanying \texttt{check\_contacts.py} constructs the cotangent
Taylor jet from Bernoulli numbers and independently checks it against
the sine and cosine jets.  It then verifies \eqref{contact:identity}
exactly for $M=1,\ldots,10$ and both signs of $\epsilon$, with $z$
and $\pi$ left symbolic.  Omitting the endpoint phase produces a
nonzero polynomial residual in all twenty cases.  In addition, the
script verifies the primitive's telescoping coefficients for
$k=0,\ldots,5$ with symbolic $L\ne1$, checks the $L=1$ branch, and
reproduces the boundary example $p=2$, $u=\lambda=1$: the specialized
ordinary integrand is zero, the continued spectral value is $-1$,
and the local endpoint term has limit $-1$.  All thirty-six exact
checks and all twenty deliberate corruption controls pass.  These
finite checks complement the all-index proofs; they do not replace
them.

% END sections/contact_verification.tex


\subsection{Integration and provenance artifacts}

\texttt{INTEGRATION.md} identifies the source questions and the
proposed scoped status updates. \texttt{claims.json} lists the
principal statements, domains, proof labels, dependencies, and
remaining limitations. \texttt{provenance/source\_manifest.json}
records the pinned source revision and hashes of inspected
manuscript files and incoming archives.
\texttt{SHA256SUMS} covers the deliverable files present at
packaging time. The complete archive is self-contained; the
original incoming ZIP files are identified by hash rather than
duplicated inside it.


% END sections/reproducibility.tex

\clearpage
\phantomsection
\addcontentsline{toc}{section}{References}
\begingroup\small
\begin{thebibliography}{99}

% BEGIN references.tex
\bibitem{ProveIt}
Vladimir Reshetnikov, \emph{ProveIt}, polylogarithm manuscript and
incoming research reports. Source revision
\snapshot, commit dated 11 October 2026, 01:57:13 UTC.
Manuscript:
\url{https://github.com/VladimirReshetnikov/ProveIt/tree/7bd45777a8a127e5dc69aff8887ca36a79a3059d/Analysis/Polylogarithms/docs/manuscript}.
Incoming reports:
\url{https://github.com/VladimirReshetnikov/ProveIt/tree/7bd45777a8a127e5dc69aff8887ca36a79a3059d/docs/incoming}.

\bibitem{IncomingCubic}
ProveIt incoming research report,
\emph{Cubic Tornheim Derivatives and Harmonic Stieltjes Constants:
A two-coordinate formula, a digamma-cube bridge, and finite reductions
at mixed integer orders}, 11 October 2026.
Archive \nolinkurl{ProveIt_Cubic_Harmonic_Mixed_Orders_2026-10-11.zip}
in the incoming directory of~\cite{ProveIt}, Sections 2 and 7,
especially further-research Question 6.

\bibitem{IncomingParity}
ProveIt incoming research report,
\emph{Harmonic Parity and Resolvent Identities:
Exact Lerch generators, Stieltjes primitives, and the quartic and higher
Tornheim layers}, 11 October 2026.
Archive
\nolinkurl{ProveIt_Harmonic_Parity_and_Resolvent_Identities_2026-10-11.zip}
in~\cite{ProveIt}, especially Sections 2, 3, and 7, Questions R1 and R4.

\bibitem{IncomingReflection}
ProveIt incoming research report,
\emph{Higher Reflection Identities}, 11 October 2026.
Archive \nolinkurl{ProveIt_Higher_Reflection_Identities_2026-10-11.zip}
in~\cite{ProveIt}, especially the harmonic section and the further
question on centered products beyond a trigamma square.

\bibitem{IncomingOrdered}
ProveIt incoming research report,
\emph{Ordered Hurwitz Germs at Every Depth:
Harmonic renormalization, directional Stieltjes constants, and
polylogarithmic Mellin identities}, 10 October 2026.
Archive \nolinkurl{ProveIt_Ordered_Hurwitz_Germs_2026-10-10.zip}
in~\cite{ProveIt}.

\bibitem{IncomingExact}
ProveIt incoming research report,
\emph{Exact Resonant Identities Beyond the Diagonal},
11 October 2026.
Archive \nolinkurl{ProveIt_Exact_Resonant_Identities_2026-10-11.zip}
in~\cite{ProveIt}, Section 5 and the Cayley verification artifacts.

\bibitem{IncomingIndependent}
ProveIt incoming research report,
\emph{Independent Orders and Exact Reductions in
Polylogarithm--Zeta Calculus}, 11 October 2026.
Archive \nolinkurl{ProveIt_Independent_Orders_2026-10-11.zip}
in~\cite{ProveIt}, Section 6 and
\nolinkurl{results/s6_incomplete_search.json}.

\bibitem{DLMF}
F. W. J. Olver et al., eds.,
\emph{NIST Digital Library of Mathematical Functions}.
Relevant formulas: hyperbolic partial fractions (4.36.5);
Gamma and digamma recurrence, reflection, Taylor expansions,
and asymptotics (Sections 5.5, 5.7, 5.11);
the Binet integral (5.9.15); the regularized Gauss integral
(15.6.1); Bernoulli polynomials (24.2);
Hurwitz zeta (25.11); polylogarithms (25.12);
Euler sums (25.16.8 and 25.16.13).
\url{https://dlmf.nist.gov/}. Consulted 11 October 2026.

\bibitem{BBG}
D. Borwein, J. M. Borwein, and R. Girgensohn,
\emph{Explicit evaluation of Euler sums},
Proceedings of the Edinburgh Mathematical Society \textbf{38}
(1995), 277--294.
\url{https://doi.org/10.1017/S0013091500019088}.
The odd-weight evaluation is displayed on p.~278 and proved in
Section 3; the boundary formula is discussed separately.

\bibitem{EspinosaMoll}
O. Espinosa and V. H. Moll,
\emph{On some integrals involving the Hurwitz zeta function:
Part 1}, The Ramanujan Journal \textbf{6} (2002), 159--188.
\url{https://www.math.tulane.edu/~vhm/papers_html/hurwitz1.pdf}.

\bibitem{Coppo}
Marc-Antoine Coppo,
\emph{On some remarkable identities related to the harmonic zeta
function}, Research in Number Theory \textbf{11} (2025).
Author manuscript:
\url{https://math.univ-cotedazur.fr/u/coppo/RNTB.pdf}.


% END references.tex

\end{thebibliography}
\endgroup
\end{document}

